% Luc Destombe --- licence CC BY-NC-SA 4.0
% https://creativecommons.org/licenses/by-nc-sa/4.0/deed.fr
% Fichier autonome : compiler avec pdflatex, deux fois (signets, nombre de pages).
\documentclass[11pt,a4paper]{article}
\makeatletter
\newif\iflycee@unecolonne
\lycee@unecolonnetrue
\newif\iflycee@palatino
\lycee@palatinotrue

\RequirePackage[utf8]{inputenc}
\RequirePackage[T1]{fontenc}
\RequirePackage[french]{babel}
\RequirePackage{etoolbox}

\newcommand{\ShorthandsOff}{\@ifundefined{shorthandoff}{}{\shorthandoff{;:!?}}}
\newcommand{\ShorthandsOn}{\@ifundefined{shorthandon}{}{\shorthandon{;:!?}}}
\ShorthandsOff
\AtBeginDocument{\ShorthandsOn}
\AtBeginEnvironment{tikzpicture}{\ShorthandsOff}

\RequirePackage{geometry}
\RequirePackage{fancyhdr}
\RequirePackage{lastpage}
\RequirePackage{multicol}
\RequirePackage{array}
\RequirePackage{enumitem}
\RequirePackage{xcolor}
\RequirePackage{needspace}

\RequirePackage{amsmath,amssymb}
\RequirePackage{mathpazo}
\DeclareMathSizes{10.95}{10.95}{8}{5.5}
\begingroup\catcode`\_=\active \catcode`\^=\active
\gdef_#1{\sb{\scriptscriptstyle #1}}
\gdef^#1{\sp{\scriptscriptstyle\lycee@moinsepais #1}}
\endgroup
\newcommand\lycee@moins{\mathbin{%
  \setbox\z@\hbox{$\m@th\scriptscriptstyle\lycee@moinsorig$}%
  \dimen@=.55\wd\z@ \advance\dimen@-.5pt
  \hbox to.75\wd\z@{\hss$\m@th\scriptscriptstyle\vcenter{\hbox{%
    \tikz\draw[line width=.5pt,line cap=round](0,0)--(\the\dimen@,0);}}$\hss}}}
\begingroup\lccode`\~=`\- \lowercase{\endgroup\let~}\lycee@moins
\newcommand\lycee@moinsepais{\mathcode`\-="8000 }
\AtBeginDocument{\mathchardef\lycee@moinsorig=\mathcode`\-\relax}
\AtBeginDocument{\catcode`\_=12 \mathcode`\_="8000
  \catcode`\^=12 \mathcode`\^="8000 }
\newcommand\lycee@exposants{%
  \fontdimen13\textfont2=.48\fontdimen6\textfont2
  \fontdimen14\textfont2=.48\fontdimen6\textfont2
  \fontdimen15\textfont2=.48\fontdimen6\textfont2
  \fontdimen13\scriptfont2=.48\fontdimen6\scriptfont2
  \fontdimen14\scriptfont2=.48\fontdimen6\scriptfont2
  \fontdimen15\scriptfont2=.48\fontdimen6\scriptfont2}
\every@math@size\expandafter{\the\every@math@size\lycee@exposants}
\newlength{\retraitcalculs}
\setlength{\retraitcalculs}{2em}

\RequirePackage{tikz}
\usetikzlibrary{arrows.meta,calc,babel,shapes.misc}
\RequirePackage[most]{tcolorbox}
\tcbuselibrary{skins,breakable}

\definecolor{coulDef}{HTML}{1F4E79}
\definecolor{coulProp}{HTML}{7030A0}
\definecolor{coulRem}{HTML}{C55A11}
\definecolor{coulEx}{HTML}{2E7D32}
\definecolor{coulExo}{HTML}{B03A2E}
\definecolor{coulPart}{HTML}{1F4E79}
\definecolor{coulMeth}{HTML}{1B6E4F}
\definecolor{coulCourbe}{HTML}{0B5ED7}
\definecolor{coulCourbeB}{HTML}{C00000}
\definecolor{coulCourbeC}{HTML}{2E7D32}
\definecolor{coulGrille}{HTML}{8C8C8C}
\definecolor{coulRep}{HTML}{1B6E4F}
\definecolor{coulBar}{HTML}{2E7D32}

\newcommand{\R}{\mathbb{R}}
\newcommand{\Rs}{\mathbb{R}^{*}}
\renewcommand{\le}{\leqslant}
\renewcommand{\ge}{\geqslant}
\renewcommand{\approx}{\simeq}
\newcommand{\vect}[1]{\overrightarrow{#1}}
\newcommand{\norme}[1]{\left\lVert #1 \right\rVert}
\newcommand{\intervalle}[2]{\left[#1\,;#2\right]}
\RequirePackage{cancel}
\renewcommand{\CancelColor}{\color{coulExo}}
\newcommand{\fois}[1]{\times{\color{coulExo}#1}}
\newcommand{\barre}[1]{\cancel{#1}}

\tikzset{grille/.style={coulGrille,line width=0.4pt}}
\newcommand{\quadrillage}[4]{\draw[grille,step=1] (#1,#2) grid (#3,#4);}
\tikzset{
  croix/.style={cross out,draw,line width=0.6pt,inner sep=0pt,minimum size=4pt},
  croixdroite/.style={croix,rotate=45,line width=0.8pt},
}
\newcommand{\pt}[3]{\path (#1) node[croix]{} node[#2,font=\small,inner sep=2.5pt]{$#3$};}

\newcommand{\lycee@repere}[4]{%
  \draw[grille,step=1] (#1,#2) grid (#3,#4);
  \draw[-{Stealth[length=2mm]},thick] (#1,0)--({#3+0.4},0) node[below left=-1pt]{$x$};
  \draw[-{Stealth[length=2mm]},thick] (0,#2)--(0,{#4+0.4}) node[below left=-1pt]{$y$};
  \node[below left,font=\scriptsize,inner sep=1.5pt] at (0,0) {$0$};}
\newcommand{\lycee@gradx}[1]{\foreach \k/\l in {#1}{%
  \draw (\k,0.07)--(\k,-0.07) node[below,font=\scriptsize,inner sep=1.5pt]{$\l$};}}
\newcommand{\lycee@grady}[1]{\foreach \k/\l in {#1}{%
  \draw (0.07,\k)--(-0.07,\k) node[left,font=\scriptsize,inner sep=1.5pt]{$\l$};}}
\newcommand{\lycee@intff}[2]{\left[#1\,;#2\right]}
\newcommand{\lycee@intfo}[2]{\left[#1\,;#2\right[}
\newcommand{\lycee@intof}[2]{\left]#1\,;#2\right]}
\newcommand{\lycee@intoo}[2]{\left]#1\,;#2\right[}
\newcommand{\lycee@ens}[1]{\left\{#1\right\}}
\AtBeginDocument{%
  \providecommand{\repere}{\lycee@repere}%
  \providecommand{\gradx}{\lycee@gradx}%
  \providecommand{\grady}{\lycee@grady}%
  \providecommand{\Cf}{\mathcal{C}\sb{\scriptscriptstyle f}}%
  \providecommand{\Cg}{\mathcal{C}\sb{\scriptscriptstyle g}}%
  \providecommand{\intff}{\lycee@intff}%
  \providecommand{\intfo}{\lycee@intfo}%
  \providecommand{\intof}{\lycee@intof}%
  \providecommand{\intoo}{\lycee@intoo}%
  \providecommand{\ens}{\lycee@ens}}

\newlist{questions}{enumerate}{3}
\setlist[questions,1]{label=\arabic*\textdegree),leftmargin=*,itemsep=2pt,topsep=3pt}
\setlist[questions,2]{label=\alph*),leftmargin=*,itemsep=1pt,topsep=2pt}
\setlist[questions,3]{label=\roman*),leftmargin=*,itemsep=1pt,topsep=1pt}
\setlist[itemize]{label=\textbullet,leftmargin=*,itemsep=2pt,topsep=3pt}

\newcommand{\besoinplace}[1]{\par\penalty\@M\begingroup
  \dimen@=#1\relax\dimen@ii=\pagegoal\advance\dimen@ii-\pagetotal
  \ifdim\dimen@>\dimen@ii\ifdim\dimen@ii>\z@\vfil\fi\break\fi\endgroup}

\newcommand{\lycee@pied}{\thepage/\pageref{LastPage}}
\newcommand{\entete}[2]{%
  \pagestyle{fancy}\fancyhf{}%
  \fancyhead[L]{\footnotesize\color{coulPart}\textbf{#1}}%
  \fancyhead[R]{\footnotesize\color{coulPart}\textbf{#2}}%
  \fancyfoot[C]{\footnotesize\lycee@pied}%
  \renewcommand{\headrulewidth}{0.4pt}%
  \renewcommand{\headrule}{\hbox to\headwidth{%
    \color{coulPart!40}\leaders\hrule height \headrulewidth\hfill}}}

\RequirePackage{ccicons}
\newcommand{\lycee@licence}{{\scriptsize\color{gray}%
  \copyright\ Luc Destombe --- \ccbyncsa\ CC BY-NC-SA 4.0}}
\newif\iflycee@licence \lycee@licencetrue
\newcommand{\sanslicence}{\lycee@licencefalse}
\AtBeginDocument{\iflycee@licence\AddToHookNext{shipout/foreground}{%
  \put(0,0){\rlap{\kern\dimexpr1in+\hoffset+\oddsidemargin+\textwidth\relax
    \llap{\raisebox{-\dimexpr1in+\voffset+\topmargin+\headheight+\headsep
      +\textheight+\footskip\relax}{\lycee@licence}}}}}\fi}

\newcommand{\formulesserrees}{%
  \setlength{\abovedisplayskip}{3pt}\setlength{\belowdisplayskip}{3pt}%
  \setlength{\abovedisplayshortskip}{2pt}\setlength{\belowdisplayshortskip}{3pt}}

\setlength{\parindent}{0pt}

\RequirePackage[implicit=false,hidelinks,pdfencoding=auto,psdextra,bookmarksopen,
  bookmarksopenlevel=1,pdfcreator={},pdfproducer={}]{hyperref}
\RequirePackage{bookmark}
\hypersetup{pdfauthor={Luc Destombe},pdfkeywords={CC BY-NC-SA 4.0}}
\newcounter{lycee@signet}
\newcommand{\signet}[3][]{\stepcounter{lycee@signet}%
  \edef\lycee@dest{\if\relax\detokenize{#1}\relax
    signet.\arabic{lycee@signet}\else#1\fi}%
  \hypertarget{\lycee@dest}{}\bookmark[level=#2,dest=\lycee@dest]{#3}}
\pdfstringdefDisableCommands{%
  \def\R{R}\def\vect#1{#1}\def\Cf{Cf}\def\Cg{Cg}\def\fois#1{×#1}%
  \def\barre#1{#1}\def\intervalle#1#2{[#1;#2]}\def\intff#1#2{[#1;#2]}%
  \def\textsuperscript#1{#1}\def\dfrac#1#2{#1/#2}\def\frac#1#2{#1/#2}%
  \def\vec#1{#1⃗}}%
\RequirePackage[official]{eurosym}

\tcbset{exercice corrige/.style={
  enhanced,breakable,before upper=\raggedright,
  colback=white,colframe=coulExo,
  boxrule=0.8pt,arc=2pt,
  left=6pt,right=6pt,top=8pt,bottom=5pt,
  before skip=9pt,after skip=4pt,
  coltitle=white,fonttitle=\bfseries\small,
  attach boxed title to top left={xshift=8pt,yshift=-\tcboxedtitleheight/2},
  boxed title style={colback=coulExo,arc=2pt,boxrule=0pt}}}
\newcounter{exo}
\newtcolorbox{exercice}[1][]{exercice corrige,
  phantom={\signet[exercice-\arabic{exo}]{2}{Exercice \arabic{exo}}},
  title={Exercice \theexo\ifx\relax#1\relax\else{} --- #1\fi}}
\newenvironment{exo}[1][\relax]{\refstepcounter{exo}\begin{exercice}[#1]}{\end{exercice}}

\geometry{top=1.8cm,bottom=1cm,left=1cm,right=1cm,headheight=14pt,footskip=0.6cm}
\renewcommand{\lycee@pied}{\thepage}

\renewcommand{\theexo}{\ifnum\value{exo}<10 0\fi\arabic{exo}}
\tcbset{exercice corrige/.style={
  enhanced,breakable,before upper=\raggedright,
  colback=white,colframe=coulExo,
  boxrule=0.9pt,arc=2pt,
  left=8pt,right=8pt,top=8pt,bottom=6pt,
  before skip=8pt,after skip=6pt,
  coltitle=white,fonttitle=\bfseries,
  attach boxed title to top left={xshift=8pt,yshift=-\tcboxedtitleheight/2},
  boxed title style={colback=coulExo,arc=2pt,boxrule=0pt}}}

\newcommand{\entetefiche}[2]{%
  \begin{center}
    \begin{tcolorbox}[enhanced,width=0.92\linewidth,colback=white,colframe=coulPart,
        boxrule=1.4pt,arc=3pt,halign=center,top=6pt,bottom=6pt]
      {\large\bfseries\color{coulPart} CORRIGÉ --- FICHE D'EXERCICES}\\[3pt]
      {\normalsize\bfseries #1}\\[2pt]
      {\normalsize #2}
    \end{tcolorbox}
  \end{center}}

\newcounter{partie}
\newcommand{\partie}[1]{%
  \stepcounter{partie}%
  \besoinplace{8\baselineskip}\vspace{4pt}%
  \begin{tcolorbox}[enhanced,colback=coulPart!8,colframe=coulPart,boxrule=0pt,
      leftrule=3pt,arc=0pt,left=6pt,right=4pt,top=3pt,bottom=3pt,
      before skip=6pt,after skip=2pt]
    \bfseries\color{coulPart}\Roman{partie} --- #1\signet{1}{\Roman{partie} --- #1}
  \end{tcolorbox}}
\newcommand{\partiesansnum}[1]{%
  \besoinplace{8\baselineskip}\vspace{4pt}%
  \begin{tcolorbox}[enhanced,colback=coulPart!8,colframe=coulPart,boxrule=0pt,
      leftrule=3pt,arc=0pt,left=6pt,right=4pt,top=3pt,bottom=3pt,
      before skip=6pt,after skip=2pt]
    \bfseries\color{coulPart}#1\signet{1}{#1}
  \end{tcolorbox}}

\setlist[questions,1]{label=\arabic*\textdegree),leftmargin=*,itemsep=3pt,topsep=3pt}
\setlist[questions,2]{label=\alph*),leftmargin=*,itemsep=2pt,topsep=2pt}
\setlist[questions,3]{label=\roman*),leftmargin=*,itemsep=1pt,topsep=1pt}
\newlist{expr}{enumerate}{1}
\setlist[expr]{label={\Alph*\ $=$},leftmargin=*,itemsep=2pt,topsep=3pt}

\newcommand{\res}[1]{{\color{coulRep}#1}}
\newcommand{\rem}[1]{\par\smallskip{\small\itshape\color{black!60}$\rightarrow$ #1}}
\newenvironment{calculs}{\par\smallskip\noindent$\displaystyle\begin{aligned}[t]}%
  {\end{aligned}$\par\smallskip}
\newcommand{\justif}[1]{\quad\text{\small\itshape\color{black!60}#1}}
\newcommand{\cas}[1]{\par\smallskip\noindent\textbf{\color{coulExo}#1}\nobreak}
\newcommand{\explique}[2]{\par\smallskip\noindent
  \begin{minipage}[t]{0.6\linewidth}\raggedright #1\end{minipage}\hfill
  \begin{minipage}[t]{0.37\linewidth}\raggedright\leavevmode
    {\small\itshape\color{black!60}#2}\end{minipage}\par}

\setlength{\parskip}{2pt}

\makeatother
\usepackage{tkz-tab}
\entete{Chapitre 1 --- Calcul littéral --- corrigé}{1\textsuperscript{re} STI2D}
\setlength{\parindent}{0pt}
\setlength{\jot}{5pt}

\let\casligne\cas
\renewcommand{\cas}[1]{\casligne{#1}\par\nopagebreak}

\newenvironment{tabsignes}[3][3.4]{\par\smallskip\begin{center}\begin{tikzpicture}
  \tkzTabInit[lgt=#1,espcl=2.4]{#2}{#3}}{\end{tikzpicture}\end{center}}

\begin{document}

\entetefiche{Chapitre 1 --- Calcul littéral}{Première STI2D}

\partie{Réduire une expression littérale}

\begin{exo}[Réduire une somme]
\explique{Pour réduire une somme, on rapproche les termes semblables (les $a$ ensemble,
les $a^{2}$ ensemble, les nombres ensemble), chacun avec le signe écrit devant lui, puis
on additionne leurs coefficients.}{définition 1 et méthode 1 du cours}
\cas{a)\quad A}
\begin{calculs}
A &= 4a+3+7a &&\justif{on recopie}\\
A &= 4a+7a+3 &&\justif{on rapproche les termes en $a$}\\
A &= \res{11a+3} &&\justif{$4+7=11$ ; le $3$ n'a pas de terme semblable}
\end{calculs}
\cas{B}
\begin{calculs}
B &= 9a+2a+5a &&\justif{trois termes semblables}\\
B &= (9+2+5)a &&\justif{on additionne les coefficients}\\
B &= \res{16a} &&
\end{calculs}
\cas{C}
\begin{calculs}
C &= 6+2a+8a &&\justif{on recopie}\\
C &= \res{10a+6} &&\justif{$2a+8a=10a$ ; le $6$ reste seul}
\end{calculs}
\rem{Erreur fréquente : écrire $C=16a$. Le nombre $6$ et les termes en $a$ ne sont pas
semblables : $10a+6$ ne se réduit plus.}
\cas{D}
\begin{calculs}
D &= 3a^{2}+7a^{2}+5a^{2} &&\justif{trois termes en $a^{2}$}\\
D &= (3+7+5)a^{2} &&\\
D &= \res{15a^{2}} &&
\end{calculs}
\cas{E}
\begin{calculs}
E &= 2a^{2}+9+4a^{2} &&\justif{on recopie}\\
E &= 2a^{2}+4a^{2}+9 &&\justif{on rapproche les $a^{2}$}\\
E &= \res{6a^{2}+9} &&
\end{calculs}
\cas{F}
\begin{calculs}
F &= 3a+6a^{2}+2a^{2} &&\justif{on recopie}\\
F &= 6a^{2}+2a^{2}+3a &&\justif{les $a^{2}$ ensemble ; $3a$ n'est pas semblable à $a^{2}$}\\
F &= \res{8a^{2}+3a} &&\justif{on s'arrête : $8a^{2}$ et $3a$ ne se regroupent pas}
\end{calculs}
\cas{b)\quad G}
\begin{calculs}
G &= 7m-12m+3m &&\justif{on recopie}\\
G &= (7-12+3)m &&\justif{de gauche à droite : $7-12=-5$, puis $-5+3=-2$}\\
G &= \res{-2m} &&
\end{calculs}
\cas{H}
\begin{calculs}
H &= 6m-9m+4m &&\\
H &= (6-9+4)m &&\justif{$6-9=-3$, puis $-3+4=1$}\\
H &= \res{m} &&\justif{$1m$ s'écrit $m$}
\end{calculs}
\cas{I}
\begin{calculs}
I &= -5m-4m-6m &&\justif{trois termes négatifs}\\
I &= (-5-4-6)m &&\justif{on ajoute les distances à zéro : $5+4+6=15$}\\
I &= \res{-15m} &&\justif{on garde le signe moins}
\end{calculs}
\cas{J}
\begin{calculs}
J &= 9m^{2}+2m^{2}-7m^{2} &&\\
J &= (9+2-7)m^{2} &&\\
J &= \res{4m^{2}} &&
\end{calculs}
\cas{K}
\begin{calculs}
K &= 5m-8m^{2}+3m^{2} &&\justif{on recopie}\\
K &= -8m^{2}+3m^{2}+5m &&\justif{les $m^{2}$ avec leur signe, puis les $m$}\\
K &= \res{-5m^{2}+5m} &&\justif{$-8+3=-5$}
\end{calculs}
\cas{L}
\begin{calculs}
L &= -4m^{2}+6m^{2}-9m^{2} &&\\
L &= (-4+6-9)m^{2} &&\justif{$-4+6=2$, puis $2-9=-7$}\\
L &= \res{-7m^{2}} &&
\end{calculs}
\cas{c)\quad M}
\begin{calculs}
M &= 3n^{2}+4n-6+5n^{2}+1+8 &&\justif{on recopie}\\
M &= 3n^{2}+5n^{2}+4n-6+1+8 &&\justif{les $n^{2}$, puis les $n$, puis les nombres}\\
M &= \res{8n^{2}+4n+3} &&\justif{$-6+1+8=3$}
\end{calculs}
\cas{N}
\begin{calculs}
N &= -7n^{2}+9-5n+3n^{2}-n-4 &&\justif{on recopie}\\
N &= -7n^{2}+3n^{2}-5n-n+9-4 &&\justif{chaque terme garde le signe écrit devant lui}\\
N &= \res{-4n^{2}-6n+5} &&\justif{$-7+3=-4$ ; $-5-1=-6$ ; $9-4=5$}
\end{calculs}
\cas{O}
\begin{calculs}
O &= -6n^{2}-3+8-2n^{2}+5n-9n &&\justif{on recopie}\\
O &= -6n^{2}-2n^{2}+5n-9n-3+8 &&\justif{on rapproche les termes semblables}\\
O &= \res{-8n^{2}-4n+5} &&\justif{$-6-2=-8$ ; $5-9=-4$ ; $-3+8=5$}
\end{calculs}
\cas{P}
\begin{calculs}
P &= 4+6n-3n^{2}+7n^{2}+2n-4 &&\justif{on recopie}\\
P &= -3n^{2}+7n^{2}+6n+2n+4-4 &&\justif{on rapproche les termes semblables}\\
P &= \res{4n^{2}+8n} &&\justif{$4-4=0$ : il ne reste pas de nombre seul}
\end{calculs}
\rem{Contrôle rapide : on remplace la lettre par $1$ dans l'expression de départ et dans le
résultat. Pour $N$ : $-7+9-5+3-1-4=-5$ et $-4-6+5=-5$.}
\end{exo}

\begin{exo}[Réduire un produit]
\explique{Pour réduire un produit, on multiplie les nombres entre eux (règle des signes),
puis les lettres entre elles : $x\times x=x^{2}$.}{méthode 1 du cours : un produit se
réduit toujours}
\cas{a)\quad A}
\begin{calculs}
A &= 8\times 2x &&\justif{on recopie}\\
A &= 8\times 2\times x &&\justif{les nombres, puis la lettre}\\
A &= \res{16x} &&
\end{calculs}
\cas{B}
\begin{calculs}
B &= 3x\times 4x &&\justif{on recopie}\\
B &= 3\times 4\times x\times x &&\justif{les nombres, puis les lettres}\\
B &= \res{12x^{2}} &&\justif{$x\times x=x^{2}$}
\end{calculs}
\cas{C}
\begin{calculs}
C &= 5\times 2x^{2} &&\\
C &= \res{10x^{2}} &&\justif{$5\times 2=10$ ; le $x^{2}$ ne change pas}
\end{calculs}
\cas{D}
\begin{calculs}
D &= 7x\times 3 &&\\
D &= \res{21x} &&\justif{une seule lettre : pas de carré}
\end{calculs}
\cas{E}
\begin{calculs}
E &= 3x^{2}\times 6 &&\\
E &= \res{18x^{2}} &&\justif{$3\times 6=18$}
\end{calculs}
\cas{F}
\begin{calculs}
F &= 4x\times 5x &&\\
F &= 4\times 5\times x\times x &&\\
F &= \res{20x^{2}} &&
\end{calculs}
\cas{b)\quad G}
\begin{calculs}
G &= -2x\times 6x &&\justif{signes contraires : le produit est négatif}\\
G &= -2\times 6\times x\times x &&\\
G &= \res{-12x^{2}} &&
\end{calculs}
\cas{H}
\begin{calculs}
H &= -9\times 3x &&\justif{signes contraires : négatif}\\
H &= \res{-27x} &&
\end{calculs}
\cas{I}
\begin{calculs}
I &= 5x\times(-2x) &&\justif{signes contraires : négatif}\\
I &= 5\times(-2)\times x\times x &&\\
I &= \res{-10x^{2}} &&
\end{calculs}
\cas{J}
\begin{calculs}
J &= -3x\times(-7x) &&\justif{deux nombres négatifs : le produit est positif}\\
J &= (-3)\times(-7)\times x\times x &&\\
J &= \res{21x^{2}} &&
\end{calculs}
\cas{K}
\begin{calculs}
K &= 4\times\left(-3x^{2}\right) &&\justif{signes contraires : négatif}\\
K &= \res{-12x^{2}} &&
\end{calculs}
\cas{L}
\begin{calculs}
L &= -5x^{2}\times 3 &&\justif{signes contraires : négatif}\\
L &= \res{-15x^{2}} &&
\end{calculs}
\cas{M}
\begin{calculs}
M &= -6\times\left(-2x^{2}\right) &&\justif{même signe : positif}\\
M &= \res{12x^{2}} &&
\end{calculs}
\cas{N}
\begin{calculs}
N &= 3x^{2}\times(-4) &&\justif{signes contraires : négatif}\\
N &= \res{-12x^{2}} &&
\end{calculs}
\cas{O}
\begin{calculs}
O &= -7x\times 2x &&\justif{signes contraires : négatif}\\
O &= -7\times 2\times x\times x &&\\
O &= \res{-14x^{2}} &&
\end{calculs}
\rem{Ne pas confondre somme et produit : $3x+4x=7x$, mais $3x\times 4x=12x^{2}$.}
\end{exo}

\begin{exo}[Somme ou produit ?]
\explique{On regarde d'abord l'opération : un produit se réduit toujours ; une somme
seulement si ses termes sont semblables.}{remarque de la partie I}
\cas{A}
\explique{$A=5+4x$ est une somme d'un nombre et d'un terme en $x$ : ils ne sont pas
semblables, $A$ \res{ne se réduit pas}.}{$5+4x$ n'est pas $9x$}
\cas{B}
\begin{calculs}
B &= 5\times 4x &&\justif{un produit}\\
B &= \res{20x} &&\justif{$5\times 4=20$}
\end{calculs}
\cas{C}
\begin{calculs}
C &= 5x\times 4x &&\justif{un produit}\\
C &= 5\times 4\times x\times x &&\\
C &= \res{20x^{2}} &&
\end{calculs}
\cas{D}
\begin{calculs}
D &= 5x+4x &&\justif{une somme de deux termes semblables}\\
D &= (5+4)x &&\justif{on additionne les coefficients}\\
D &= \res{9x} &&
\end{calculs}
\explique{Seule $A$ ne se réduit pas : c'est une somme de termes qui ne sont pas
semblables.}{réponse à la question}
\rem{Erreur fréquente : $5x+4x=20x^{2}$. Dans une somme, on additionne les coefficients et
la lettre ne change pas.}
\end{exo}

\partie{Développer}

\begin{exo}[Distributivité simple]
\explique{Le facteur devant (ou derrière) la parenthèse multiplie \emph{chaque} terme :
$k(a+b)=ka+kb$. On réduit ensuite.}{propriété 1 et méthode 2 du cours}
\cas{a)\quad A}
\begin{calculs}
A &= 4(3x+5) &&\justif{on recopie}\\
A &= 4\times 3x+4\times 5 &&\justif{$4$ multiplie $3x$, puis $5$}\\
A &= \res{12x+20} &&
\end{calculs}
\cas{B}
\begin{calculs}
B &= (6+5x)\times 2 &&\justif{le facteur est à droite : il multiplie aussi chaque terme}\\
B &= 6\times 2+5x\times 2 &&\\
B &= 12+10x &&\\
B &= \res{10x+12} &&\justif{on range : les $x$ d'abord}
\end{calculs}
\cas{C}
\begin{calculs}
C &= x(2x+9) &&\justif{le facteur est $x$}\\
C &= x\times 2x+x\times 9 &&\\
C &= \res{2x^{2}+9x} &&\justif{$x\times 2x=2x^{2}$}
\end{calculs}
\cas{D}
\begin{calculs}
D &= 3x(4x+5) &&\\
D &= 3x\times 4x+3x\times 5 &&\\
D &= \res{12x^{2}+15x} &&\justif{$3x\times 4x=12x^{2}$ et $3x\times 5=15x$}
\end{calculs}
\cas{b)\quad E}
\begin{calculs}
E &= 5x+3(2x+4) &&\justif{le produit d'abord ; le $5x$ attend}\\
E &= 5x+6x+12 &&\justif{$3\times 2x=6x$ et $3\times 4=12$}\\
E &= \res{11x+12} &&
\end{calculs}
\rem{Erreur fréquente : additionner $5x$ et $3$ avant de développer. La multiplication passe
avant l'addition.}
\cas{F}
\begin{calculs}
F &= 2(5x+6)+4(2x+3) &&\justif{deux produits, développés chacun}\\
F &= 10x+12+8x+12 &&\\
F &= \res{18x+24} &&
\end{calculs}
\cas{G}
\begin{calculs}
G &= 6(3+2x)+5x &&\\
G &= 18+12x+5x &&\\
G &= \res{17x+18} &&
\end{calculs}
\cas{H}
\begin{calculs}
H &= 3(4x+1)+2(7+3x) &&\\
H &= 12x+3+14+6x &&\\
H &= \res{18x+17} &&\justif{$12x+6x=18x$ ; $3+14=17$}
\end{calculs}
\cas{I}
\begin{calculs}
I &= x(3x+5)+2(6x+1) &&\\
I &= 3x^{2}+5x+12x+2 &&\\
I &= \res{3x^{2}+17x+2} &&
\end{calculs}
\cas{J}
\begin{calculs}
J &= 2x(4x+3)+3x(2x+5) &&\\
J &= 8x^{2}+6x+6x^{2}+15x &&\justif{$2x\times 4x=8x^{2}$ ; $3x\times 2x=6x^{2}$}\\
J &= \res{14x^{2}+21x} &&
\end{calculs}
\rem{Erreur fréquente : ne multiplier que le premier terme, $4(3x+5)=12x+5$. Contrôle avec
$x=1$ : $4\times 8=32$ et $12+20=32$.}
\end{exo}

\begin{exo}[Signes et facteurs négatifs]
\explique{Le facteur garde son signe et multiplie chaque terme : on applique la règle des
signes à chaque produit.}{propriété 1 et méthode 2 du cours, exemple a)}
\cas{a)\quad A}
\begin{calculs}
A &= 5(a+3)-4a &&\justif{on recopie}\\
A &= 5a+15-4a &&\\
A &= \res{a+15} &&\justif{$5a-4a=a$}
\end{calculs}
\cas{B}
\begin{calculs}
B &= -7a+3(2a-4) &&\\
B &= -7a+6a-12 &&\justif{$3\times(-4)=-12$}\\
B &= \res{-a-12} &&\justif{$-7a+6a=-1a$, qu'on écrit $-a$}
\end{calculs}
\cas{C}
\begin{calculs}
C &= 4(2+3a)-3(5+2a) &&\justif{le second facteur est $-3$}\\
C &= 8+12a-15-6a &&\justif{$-3\times 5=-15$ et $-3\times 2a=-6a$}\\
C &= \res{6a-7} &&\justif{$12a-6a=6a$ ; $8-15=-7$}
\end{calculs}
\cas{D}
\begin{calculs}
D &= -5(-3a+2)+2(4a-3) &&\\
D &= 15a-10+8a-6 &&\justif{$-5\times(-3a)=+15a$ ; $-5\times 2=-10$}\\
D &= \res{23a-16} &&
\end{calculs}
\cas{b)\quad E}
\begin{calculs}
E &= 3b+2b(5+6b) &&\\
E &= 3b+10b+12b^{2} &&\justif{$2b\times 6b=12b^{2}$}\\
E &= \res{12b^{2}+13b} &&\justif{on range : les $b^{2}$ d'abord}
\end{calculs}
\cas{F}
\begin{calculs}
F &= 4b(-3b+5)+7b^{2} &&\\
F &= -12b^{2}+20b+7b^{2} &&\justif{$4b\times(-3b)=-12b^{2}$}\\
F &= \res{-5b^{2}+20b} &&\justif{$-12+7=-5$}
\end{calculs}
\cas{G}
\begin{calculs}
G &= 3b(-4b-2)-b(2b-9) &&\justif{le second facteur est $-b$}\\
G &= -12b^{2}-6b-2b^{2}+9b &&\justif{$-b\times 2b=-2b^{2}$ et $-b\times(-9)=+9b$}\\
G &= \res{-14b^{2}+3b} &&
\end{calculs}
\cas{H}
\begin{calculs}
H &= -2b(5b-3)+6b(-2b-4) &&\\
H &= -10b^{2}+6b-12b^{2}-24b &&\justif{$-2b\times(-3)=+6b$ ; $6b\times(-4)=-24b$}\\
H &= \res{-22b^{2}-18b} &&
\end{calculs}
\rem{Erreur fréquente dans $C$ et $G$ : ne changer que le signe du premier terme. Le
facteur $-3$ (ou $-b$) multiplie \emph{chaque} terme de la parenthèse.}
\end{exo}

\begin{exo}[Double distributivité]
\explique{Chaque terme de la première parenthèse multiplie chaque terme de la seconde :
$(a+b)(c+d)=ac+ad+bc+bd$, soit quatre produits, puis on réduit.}{propriété 1 et méthode 2
du cours, exemple b)}
\cas{A}
\begin{calculs}
A &= (2x+5)(4x+3) &&\justif{on recopie}\\
A &= 8x^{2}+6x+20x+15 &&\justif{$2x\times 4x$ ; $2x\times 3$ ; $5\times 4x$ ; $5\times 3$}\\
A &= \res{8x^{2}+26x+15} &&
\end{calculs}
\cas{B}
\begin{calculs}
B &= (3+4x)(5x+2) &&\\
B &= 15x+6+20x^{2}+8x &&\justif{$3\times 5x$ ; $3\times 2$ ; $4x\times 5x$ ; $4x\times 2$}\\
B &= \res{20x^{2}+23x+6} &&\justif{on réduit et on range}
\end{calculs}
\cas{C}
\begin{calculs}
C &= (5x+2)(4+3x) &&\\
C &= 20x+15x^{2}+8+6x &&\\
C &= \res{15x^{2}+26x+8} &&
\end{calculs}
\cas{D}
\begin{calculs}
D &= (6x+1)(3+5x) &&\\
D &= 18x+30x^{2}+3+5x &&\\
D &= \res{30x^{2}+23x+3} &&
\end{calculs}
\cas{E}
\begin{calculs}
E &= 6x+5(2+3x) &&\justif{un seul produit : $5$ multiplie $(2+3x)$}\\
E &= 6x+10+15x &&\\
E &= \res{21x+10} &&
\end{calculs}
\cas{F}
\begin{calculs}
F &= (4x+3)(4x+3) &&\\
F &= 16x^{2}+12x+12x+9 &&\\
F &= \res{16x^{2}+24x+9} &&
\end{calculs}
\explique{$E$ n'est pas un produit de deux parenthèses : seul le nombre $5$ multiplie
$(2+3x)$. C'est une distributivité simple, avec deux produits au lieu de quatre.}{réponse à
la question}
\rem{$F=(4x+3)^{2}$ : une identité remarquable ira plus vite (partie III).}
\end{exo}

\partie{Les trois identités remarquables}

\begin{exo}[Identité $(a+b)^{2}$]
\explique{On repère $a$ et $b$, puis on écrit $a^{2}+2ab+b^{2}$ : le carré du premier, le
double produit, le carré du second.}{propriété 2 et méthode 3 du cours}
\cas{A}
\begin{calculs}
A &= (3x+2)^{2} &&\justif{$(a+b)^{2}$, avec $a=3x$ et $b=2$}\\
A &= (3x)^{2}+2\times 3x\times 2+2^{2} &&\justif{$a^{2}+2ab+b^{2}$}\\
A &= \res{9x^{2}+12x+4} &&\justif{$(3x)^{2}=9x^{2}$}
\end{calculs}
\rem{$(3x)^{2}=9x^{2}$, et non $3x^{2}$ : le carré porte sur le nombre \emph{et} sur la
lettre.}
\cas{B}
\begin{calculs}
B &= (4+3x)^{2} &&\justif{$a=4$ et $b=3x$}\\
B &= 4^{2}+2\times 4\times 3x+(3x)^{2} &&\\
B &= 16+24x+9x^{2} &&\\
B &= \res{9x^{2}+24x+16} &&\justif{on range}
\end{calculs}
\cas{C}
\begin{calculs}
C &= (2a+7)^{2} &&\justif{premier terme $2a$, second terme $7$}\\
C &= (2a)^{2}+2\times 2a\times 7+7^{2} &&\\
C &= \res{4a^{2}+28a+49} &&
\end{calculs}
\rem{Ici la lettre est $a$ : on parle de premier et de second terme pour ne pas la confondre
avec le $a$ de la formule.}
\cas{D}
\begin{calculs}
D &= (5+4b)^{2} &&\justif{premier terme $5$, second terme $4b$}\\
D &= 5^{2}+2\times 5\times 4b+(4b)^{2} &&\\
D &= 25+40b+16b^{2} &&\\
D &= \res{16b^{2}+40b+25} &&
\end{calculs}
\cas{E}
\begin{calculs}
E &= (x+9)^{2} &&\justif{$a=x$ et $b=9$}\\
E &= x^{2}+2\times x\times 9+9^{2} &&\\
E &= \res{x^{2}+18x+81} &&
\end{calculs}
\cas{F}
\begin{calculs}
F &= (6x+1)^{2} &&\justif{$a=6x$ et $b=1$}\\
F &= (6x)^{2}+2\times 6x\times 1+1^{2} &&\\
F &= \res{36x^{2}+12x+1} &&
\end{calculs}
\cas{G}
\begin{calculs}
G &= (2m+5)^{2} &&\justif{$a=2m$ et $b=5$}\\
G &= (2m)^{2}+2\times 2m\times 5+5^{2} &&\\
G &= \res{4m^{2}+20m+25} &&
\end{calculs}
\cas{H}
\begin{calculs}
H &= (3+7r)^{2} &&\justif{$a=3$ et $b=7r$}\\
H &= 3^{2}+2\times 3\times 7r+(7r)^{2} &&\\
H &= 9+42r+49r^{2} &&\\
H &= \res{49r^{2}+42r+9} &&
\end{calculs}
\rem{Erreur fréquente : $(a+b)^{2}=a^{2}+b^{2}$. Le double produit $2ab$ ne s'oublie
jamais : avec $x=1$, $(x+9)^{2}=100$, alors que $1+81=82$.}
\end{exo}

\begin{exo}[Identité $(a-b)^{2}$]
\explique{On repère $a$ et $b$ (sans le signe $-$), puis on écrit $a^{2}-2ab+b^{2}$ : seul le
double produit porte le signe $-$.}{propriété 2 et méthode 3 du cours, exemple b)}
\cas{A}
\begin{calculs}
A &= (2x-3)^{2} &&\justif{$(a-b)^{2}$, avec $a=2x$ et $b=3$}\\
A &= (2x)^{2}-2\times 2x\times 3+3^{2} &&\justif{$a^{2}-2ab+b^{2}$}\\
A &= \res{4x^{2}-12x+9} &&
\end{calculs}
\cas{B}
\begin{calculs}
B &= (5x-4)^{2} &&\justif{$a=5x$ et $b=4$}\\
B &= (5x)^{2}-2\times 5x\times 4+4^{2} &&\\
B &= \res{25x^{2}-40x+16} &&
\end{calculs}
\cas{C}
\begin{calculs}
C &= (4-3x)^{2} &&\justif{$a=4$ et $b=3x$}\\
C &= 4^{2}-2\times 4\times 3x+(3x)^{2} &&\\
C &= 16-24x+9x^{2} &&\\
C &= \res{9x^{2}-24x+16} &&\justif{on range}
\end{calculs}
\cas{D}
\begin{calculs}
D &= (1-6x)^{2} &&\justif{$a=1$ et $b=6x$}\\
D &= 1^{2}-2\times 1\times 6x+(6x)^{2} &&\\
D &= 1-12x+36x^{2} &&\\
D &= \res{36x^{2}-12x+1} &&
\end{calculs}
\cas{E}
\begin{calculs}
E &= (3a-5)^{2} &&\justif{premier terme $3a$, second terme $5$}\\
E &= (3a)^{2}-2\times 3a\times 5+5^{2} &&\\
E &= \res{9a^{2}-30a+25} &&
\end{calculs}
\cas{F}
\begin{calculs}
F &= (2-7b)^{2} &&\justif{premier terme $2$, second terme $7b$}\\
F &= 2^{2}-2\times 2\times 7b+(7b)^{2} &&\\
F &= 4-28b+49b^{2} &&\\
F &= \res{49b^{2}-28b+4} &&
\end{calculs}
\cas{G}
\begin{calculs}
G &= (4m-1)^{2} &&\justif{$a=4m$ et $b=1$}\\
G &= (4m)^{2}-2\times 4m\times 1+1^{2} &&\\
G &= \res{16m^{2}-8m+1} &&
\end{calculs}
\cas{H}
\begin{calculs}
H &= (n-8)^{2} &&\justif{$a=n$ et $b=8$}\\
H &= n^{2}-2\times n\times 8+8^{2} &&\\
H &= \res{n^{2}-16n+64} &&
\end{calculs}
\rem{Le dernier terme $b^{2}$ est toujours positif : $(n-8)^{2}$ se termine par $+64$, pas
par $-64$.}
\end{exo}

\begin{exo}[Identité $(a+b)(a-b)$]
\explique{Deux parenthèses avec les mêmes termes, l'une avec $+$, l'autre avec $-$ :
$(a+b)(a-b)=a^{2}-b^{2}$.}{propriété 2 et méthode 3 du cours, exemple c)}
\cas{A}
\begin{calculs}
A &= (3x+7)(3x-7) &&\justif{$a=3x$ et $b=7$}\\
A &= (3x)^{2}-7^{2} &&\justif{$a^{2}-b^{2}$}\\
A &= \res{9x^{2}-49} &&
\end{calculs}
\cas{B}
\explique{$(2x+3)(2x-5)$ : les seconds termes $3$ et $5$ ne sont pas les mêmes, ce
\res{n'est pas une identité remarquable}.}{on contrôle avant d'appliquer la formule}
\begin{calculs}
B &= (2x+3)(2x-5) &&\justif{double distributivité}\\
B &= 4x^{2}-10x+6x-15 &&\justif{quatre produits}\\
B &= \res{4x^{2}-4x-15} &&
\end{calculs}
\cas{C}
\begin{calculs}
C &= (5+2x)(5-2x) &&\justif{$a=5$ et $b=2x$}\\
C &= 5^{2}-(2x)^{2} &&\\
C &= \res{25-4x^{2}} &&
\end{calculs}
\cas{D}
\begin{calculs}
D &= (4a+1)(4a-1) &&\justif{premier terme $4a$, second terme $1$}\\
D &= (4a)^{2}-1^{2} &&\\
D &= \res{16a^{2}-1} &&
\end{calculs}
\cas{E}
\begin{calculs}
E &= (b-6)(b+6) &&\justif{le produit ne dépend pas de l'ordre des facteurs}\\
E &= b^{2}-6^{2} &&\justif{premier terme $b$, second terme $6$}\\
E &= \res{b^{2}-36} &&
\end{calculs}
\cas{F}
\begin{calculs}
F &= (9+c)(9-c) &&\justif{$a=9$ et $b=c$}\\
F &= 9^{2}-c^{2} &&\\
F &= \res{81-c^{2}} &&
\end{calculs}
\explique{L'expression qui n'est pas une identité remarquable est $B$.}{réponse à la
question}
\end{exo}

\begin{exo}[Choisir la bonne méthode]
\explique{Avant de calculer, on regarde la forme : un nombre devant une parenthèse
(distributivité simple), un carré $(a+b)^{2}$ ou $(a-b)^{2}$, un produit $(a+b)(a-b)$,
sinon la double distributivité.}{méthodes 2 et 3 du cours}
\cas{a)\quad A}
\begin{calculs}
A &= 3(2x+7) &&\justif{distributivité simple}\\
A &= \res{6x+21} &&\justif{$3\times 2x=6x$ et $3\times 7=21$}
\end{calculs}
\cas{B}
\begin{calculs}
B &= (2x+7)^{2} &&\justif{$(a+b)^{2}$, avec $a=2x$ et $b=7$}\\
B &= (2x)^{2}+2\times 2x\times 7+7^{2} &&\\
B &= \res{4x^{2}+28x+49} &&
\end{calculs}
\cas{C}
\begin{calculs}
C &= (5x-3)^{2} &&\justif{$(a-b)^{2}$, avec $a=5x$ et $b=3$}\\
C &= (5x)^{2}-2\times 5x\times 3+3^{2} &&\\
C &= \res{25x^{2}-30x+9} &&
\end{calculs}
\cas{D}
\begin{calculs}
D &= (7x-1)(7x+1) &&\justif{$(a-b)(a+b)$, avec $a=7x$ et $b=1$}\\
D &= (7x)^{2}-1^{2} &&\\
D &= \res{49x^{2}-1} &&
\end{calculs}
\cas{E}
\begin{calculs}
E &= (6x+5)(2x-3) &&\justif{pas d'identité : double distributivité}\\
E &= 12x^{2}-18x+10x-15 &&\\
E &= \res{12x^{2}-8x-15} &&
\end{calculs}
\cas{b)\quad F}
\begin{calculs}
F &= (3x+2)(x+4)-(2x-5)(x+1) &&\justif{on recopie}\\
F &= 3x^{2}+12x+2x+8-\left(2x^{2}+2x-5x-5\right) &&\justif{on garde la parenthèse derrière le $-$}\\
F &= 3x^{2}+14x+8-\left(2x^{2}-3x-5\right) &&\justif{on réduit dans la parenthèse}\\
F &= 3x^{2}+14x+8-2x^{2}+3x+5 &&\justif{le $-$ change tous les signes}\\
F &= \res{x^{2}+17x+13} &&\justif{on réduit et on ordonne}
\end{calculs}
\cas{G}
\begin{calculs}
G &= (x-3)^{2}+(2x+1)(2x-1) &&\justif{deux identités : $(a-b)^{2}$ et $(a+b)(a-b)$}\\
G &= x^{2}-6x+9+4x^{2}-1 &&\justif{un $+$ devant : les signes ne changent pas}\\
G &= \res{5x^{2}-6x+8} &&
\end{calculs}
\cas{H}
\begin{calculs}
H &= (2x+5)^{2}-(x-4)(3x+2) &&\justif{une identité, puis une double distributivité}\\
H &= 4x^{2}+20x+25-\left(3x^{2}+2x-12x-8\right) &&\justif{on garde la parenthèse derrière le $-$}\\
H &= 4x^{2}+20x+25-\left(3x^{2}-10x-8\right) &&\\
H &= 4x^{2}+20x+25-3x^{2}+10x+8 &&\justif{le $-$ change tous les signes}\\
H &= \res{x^{2}+30x+33} &&
\end{calculs}
\rem{Le signe $-$ devant un produit change \emph{tous} les signes : on garde la parenthèse
tant que le produit n'est pas réduit. C'est l'erreur n°1 du chapitre.}
\end{exo}

\partie{Factoriser}

\begin{exo}[Facteur commun]
\explique{On cherche le plus grand nombre qui divise tous les coefficients, et la lettre
présente dans chaque terme ; on écrit chaque terme comme un produit par ce facteur, puis
$ka+kb=k(a+b)$.}{méthode 4 du cours, exemple a)}
\cas{A}
\begin{calculs}
A &= 6a+6b &&\justif{$6$ est dans les deux termes}\\
A &= \res{6(a+b)} &&
\end{calculs}
\cas{B}
\begin{calculs}
B &= 48a-36b &&\justif{$48$ et $36$ sont dans la table de $12$}\\
B &= 12\times 4a-12\times 3b &&\\
B &= \res{12(4a-3b)} &&
\end{calculs}
\cas{C}
\begin{calculs}
C &= 45a-15b &&\\
C &= 15\times 3a-15\times b &&\justif{$45a=15\times 3a$ : le facteur commun est $15$}\\
C &= \res{15(3a-b)} &&
\end{calculs}
\cas{D}
\begin{calculs}
D &= 120a-80b &&\\
D &= 40\times 3a-40\times 2b &&\justif{$120$ et $80$ sont dans la table de $40$}\\
D &= \res{40(3a-2b)} &&
\end{calculs}
\cas{E}
\begin{calculs}
E &= 7x^{2}+4x &&\justif{$x$ est dans les deux termes}\\
E &= x\times 7x+x\times 4 &&\\
E &= \res{x(7x+4)} &&
\end{calculs}
\cas{F}
\begin{calculs}
F &= 18x^{2}-11x &&\\
F &= x\times 18x-x\times 11 &&\\
F &= \res{x(18x-11)} &&
\end{calculs}
\cas{G}
\begin{calculs}
G &= 15x+8x^{2} &&\\
G &= x\times 15+x\times 8x &&\\
G &= \res{x(15+8x)} &&
\end{calculs}
\cas{H}
\begin{calculs}
H &= 10x^{2}-4x &&\justif{$2$ et $x$ sont dans les deux termes}\\
H &= 2x\times 5x-2x\times 2 &&\\
H &= \res{2x(5x-2)} &&
\end{calculs}
\cas{I}
\begin{calculs}
I &= 4x^{2}+x &&\\
I &= x\times 4x+x\times 1 &&\justif{$x=x\times 1$}\\
I &= \res{x(4x+1)} &&
\end{calculs}
\rem{Dans $I$, il reste $1$ dans la parenthèse, pas $0$ : $x(4x)$ redonnerait seulement
$4x^{2}$.}
\cas{J}
\begin{calculs}
J &= 3x^{3}-x^{2} &&\justif{$x^{2}$ est dans les deux termes : $x^{3}=x^{2}\times x$}\\
J &= x^{2}\times 3x-x^{2}\times 1 &&\\
J &= \res{x^{2}(3x-1)} &&
\end{calculs}
\rem{Contrôle : on redéveloppe. $2x(5x-2)=10x^{2}-4x$, c'est bien $H$.}
\end{exo}

\begin{exo}[Identité $a^{2}+2ab+b^{2}$]
\explique{Trois termes dont deux carrés : on trouve $a$ et $b$, puis on contrôle que le
troisième terme est bien le double produit $2ab$. Si oui, l'expression est $(a+b)^{2}$.}{méthode
4 du cours, exemple d)}
\cas{A}
\begin{calculs}
A &= 9x^{2}+12x+4 &&\justif{$9x^{2}=(3x)^{2}$ et $4=2^{2}$}\\
A &= (3x)^{2}+2\times 3x\times 2+2^{2} &&\justif{le double produit vaut bien $12x$}\\
A &= \res{(3x+2)^{2}} &&
\end{calculs}
\cas{B}
\begin{calculs}
B &= 25x^{2}+40x+16 &&\justif{$25x^{2}=(5x)^{2}$ et $16=4^{2}$}\\
B &= (5x)^{2}+2\times 5x\times 4+4^{2} &&\justif{$2\times 5x\times 4=40x$ : c'est bien le terme du milieu}\\
B &= \res{(5x+4)^{2}} &&
\end{calculs}
\cas{C}
\explique{$9a^{2}=(3a)^{2}$ et $4=2^{2}$ ; le double produit serait $2\times 3a\times 2$, soit
$12a$. Or le terme du milieu est $6a$ : $C$ \res{n'est pas un carré}.}{le contrôle échoue}
\cas{D}
\begin{calculs}
D &= 25+30b+9b^{2} &&\justif{$25=5^{2}$ et $9b^{2}=(3b)^{2}$}\\
D &= 5^{2}+2\times 5\times 3b+(3b)^{2} &&\justif{$2\times 5\times 3b=30b$}\\
D &= \res{(5+3b)^{2}} &&
\end{calculs}
\cas{E}
\explique{$16c^{2}=(4c)^{2}$ et $25=5^{2}$ ; le double produit serait $2\times 4c\times 5$,
soit $40c$. Or le terme du milieu est $20c$ : $E$ \res{n'est pas un carré}.}{le contrôle
échoue}
\cas{F}
\begin{calculs}
F &= 28d+49d^{2}+4 &&\justif{on range d'abord}\\
F &= 49d^{2}+28d+4 &&\justif{$49d^{2}=(7d)^{2}$ et $4=2^{2}$}\\
F &= (7d)^{2}+2\times 7d\times 2+2^{2} &&\justif{$2\times 7d\times 2=28d$}\\
F &= \res{(7d+2)^{2}} &&
\end{calculs}
\explique{Les deux expressions qui ne sont pas des carrés sont $C$ et $E$.}{réponse à la
question}
\rem{Erreur fréquente : factoriser $9a^{2}+6a+4$ en $(3a+2)^{2}$ sans contrôler. On vérifie
en redéveloppant : $(3a+2)^{2}=9a^{2}+12a+4$.}
\end{exo}

\begin{exo}[Identités $a^{2}-2ab+b^{2}$ et $a^{2}-b^{2}$]
\explique{Trois termes : deux carrés et un double produit précédé de $-$ donnent
$(a-b)^{2}$. Deux carrés séparés par un $-$ : $a^{2}-b^{2}=(a-b)(a+b)$.}{méthode 4 du cours,
exemples d) à f)}
\cas{a)\quad A}
\begin{calculs}
A &= 16x^{2}-24x+9 &&\justif{$16x^{2}=(4x)^{2}$ et $9=3^{2}$}\\
A &= (4x)^{2}-2\times 4x\times 3+3^{2} &&\justif{$2\times 4x\times 3=24x$, précédé de $-$}\\
A &= \res{(4x-3)^{2}} &&
\end{calculs}
\cas{B}
\begin{calculs}
B &= 9-12x+4x^{2} &&\justif{$9=3^{2}$ et $4x^{2}=(2x)^{2}$}\\
B &= 3^{2}-2\times 3\times 2x+(2x)^{2} &&\justif{$2\times 3\times 2x=12x$}\\
B &= \res{(3-2x)^{2}} &&
\end{calculs}
\cas{C}
\explique{$4x^{2}=(2x)^{2}$ et $9=3^{2}$ ; le double produit serait $2\times 2x\times 3$, soit
$12x$. Or le terme du milieu est $-6x$ : $C$ \res{n'est pas un carré}.}{le contrôle échoue}
\cas{D}
\explique{$D=25x^{2}+10x-1$ : le dernier terme $-1$ est négatif, ce n'est pas un carré. $D$
\res{n'est pas un carré}.}{dans $a^{2}-2ab+b^{2}$, les deux carrés sont précédés d'un $+$}
\cas{E}
\begin{calculs}
E &= 9a^{2}+25-30a &&\justif{on range d'abord}\\
E &= 9a^{2}-30a+25 &&\justif{$9a^{2}=(3a)^{2}$ et $25=5^{2}$}\\
E &= (3a)^{2}-2\times 3a\times 5+5^{2} &&\justif{$2\times 3a\times 5=30a$}\\
E &= \res{(3a-5)^{2}} &&
\end{calculs}
\cas{F}
\begin{calculs}
F &= 49b^{2}+4-28b &&\justif{on range d'abord}\\
F &= 49b^{2}-28b+4 &&\justif{$49b^{2}=(7b)^{2}$ et $4=2^{2}$}\\
F &= (7b)^{2}-2\times 7b\times 2+2^{2} &&\justif{$2\times 7b\times 2=28b$}\\
F &= \res{(7b-2)^{2}} &&
\end{calculs}
\explique{Les deux expressions qui ne sont pas des carrés sont $C$ et $D$.}{réponse à la
question}
\cas{b)\quad G}
\begin{calculs}
G &= 49x^{2}-25 &&\justif{différence de deux carrés}\\
G &= (7x)^{2}-5^{2} &&\justif{$a^{2}-b^{2}$, avec $a=7x$ et $b=5$}\\
G &= \res{(7x-5)(7x+5)} &&
\end{calculs}
\cas{H}
\begin{calculs}
H &= 64-9x^{2} &&\\
H &= 8^{2}-(3x)^{2} &&\justif{$a=8$ et $b=3x$}\\
H &= \res{(8-3x)(8+3x)} &&
\end{calculs}
\cas{I}
\begin{calculs}
I &= (x+2)^{2}-16 &&\justif{$a^{2}-b^{2}$, avec $a=x+2$ et $b=4$}\\
I &= \left[(x+2)-4\right]\left[(x+2)+4\right] &&\justif{crochets : $a$ est une parenthèse}\\
I &= \res{(x-2)(x+6)} &&\justif{on réduit chaque crochet}
\end{calculs}
\cas{J}
\begin{calculs}
J &= 36-(x-1)^{2} &&\justif{$a^{2}-b^{2}$, avec $a=6$ et $b=x-1$}\\
J &= \left[6-(x-1)\right]\left[6+(x-1)\right] &&\justif{crochets : $b$ est une parenthèse}\\
J &= (6-x+1)(6+x-1) &&\justif{le $-$ change les signes de $x-1$}\\
J &= \res{(7-x)(5+x)} &&
\end{calculs}
\rem{Dans $J$, sans crochets, on écrit souvent $6-x-1$ : c'est l'erreur de signe à
éviter.}
\end{exo}

\begin{exo}[Facteur commun entre parenthèses]
\explique{On repère la parenthèse commune aux deux termes, on l'écrit devant, puis ce qui
reste de chaque terme dans un crochet ; on réduit le crochet.}{méthode 4 du cours,
exemples b) et c)}
\cas{A}
\begin{calculs}
A &= (3x+4)(2x-5)+(3x+4)(5x+1) &&\justif{facteur commun : $(3x+4)$}\\
A &= (3x+4)\left[(2x-5)+(5x+1)\right] &&\justif{ce qui reste, dans un crochet}\\
A &= \res{(3x+4)(7x-4)} &&\justif{$2x+5x=7x$ ; $-5+1=-4$}
\end{calculs}
\cas{B}
\begin{calculs}
B &= (6x-1)(2x+3)+(6x-1)(x-5) &&\justif{facteur commun : $(6x-1)$}\\
B &= (6x-1)\left[(2x+3)+(x-5)\right] &&\\
B &= \res{(6x-1)(3x-2)} &&
\end{calculs}
\cas{C}
\begin{calculs}
C &= (2x+9)(4x-1)+(4x-1)(3x+2) &&\justif{facteur commun : $(4x-1)$, à gauche puis à droite}\\
C &= (4x-1)\left[(2x+9)+(3x+2)\right] &&\\
C &= \res{(4x-1)(5x+11)} &&
\end{calculs}
\cas{D}
\begin{calculs}
D &= (4r-3)(2r+5)+(r+7)(2r+5) &&\justif{facteur commun : $(2r+5)$}\\
D &= (2r+5)\left[(4r-3)+(r+7)\right] &&\\
D &= \res{(2r+5)(5r+4)} &&
\end{calculs}
\cas{E}
\begin{calculs}
E &= (5x+3)(x-4)+(5x+3) &&\justif{on recopie}\\
E &= (5x+3)(x-4)+(5x+3)\times 1 &&\justif{on fait apparaître le $1$}\\
E &= (5x+3)\left[(x-4)+1\right] &&\\
E &= \res{(5x+3)(x-3)} &&
\end{calculs}
\rem{Sans le $\times 1$, on oublie souvent le $+1$ et on écrit $(5x+3)(x-4)$.}
\cas{F}
\begin{calculs}
F &= (7r-2)(3r+1)+(7r-2)^{2} &&\justif{on recopie}\\
F &= (7r-2)(3r+1)+(7r-2)(7r-2) &&\justif{le carré écrit comme un produit}\\
F &= (7r-2)\left[(3r+1)+(7r-2)\right] &&\\
F &= \res{(7r-2)(10r-1)} &&
\end{calculs}
\cas{G}
\begin{calculs}
G &= (4m-5)^{2}-(4m-5)(m+3) &&\justif{on recopie}\\
G &= (4m-5)(4m-5)-(4m-5)(m+3) &&\justif{le carré écrit comme un produit}\\
G &= (4m-5)\left[(4m-5)-(m+3)\right] &&\justif{facteur commun $(4m-5)$}\\
G &= (4m-5)(4m-5-m-3) &&\justif{le $-$ change les signes de $m+3$}\\
G &= \res{(4m-5)(3m-8)} &&
\end{calculs}
\rem{Dans $G$, le $-$ devant le second produit change les deux signes de $(m+3)$ : on
garde le crochet jusqu'à cette étape.}
\end{exo}

\partie{Résoudre des équations}

\begin{exo}[Équations du premier degré]
\explique{On regroupe les $x$ d'un côté et les nombres de l'autre, en ajoutant ou en
soustrayant le même terme aux deux membres, puis on divise par le coefficient de $x$.}{méthode
5 du cours, exemple a)}
\cas{a)}
\begin{calculs}
4x-9 &= 3x+12 &&\justif{on recopie}\\
4x-3x &= 12+9 &&\justif{les $x$ à gauche, les nombres à droite}\\
x &= 21 &&
\end{calculs}
\explique{$\res{S=\ens{21}}$.}{vérification : $84-9=75$ et $63+12=75$}
\cas{b)}
\begin{calculs}
5x-17 &= 2x-5 &&\justif{on recopie}\\
5x-2x &= -5+17 &&\\
3x &= 12 &&\\
x &= 4 &&\justif{on divise par $3$}
\end{calculs}
\explique{$\res{S=\ens{4}}$.}{vérification : $20-17=3$ et $8-5=3$}
\cas{c)}
\begin{calculs}
2x-11 &= 5x+13 &&\justif{on recopie}\\
-11-13 &= 5x-2x &&\justif{les $x$ du côté où il y en a le plus}\\
-24 &= 3x &&\\
x &= -8 &&\justif{on divise par $3$ : $-24\div 3=-8$}
\end{calculs}
\explique{$\res{S=\ens{-8}}$.}{vérification : $-16-11=-27$ et $-40+13=-27$}
\cas{d)}
\begin{calculs}
3x+8 &= 3x-4 &&\justif{on recopie}\\
3x-3x &= -4-8 &&\\
0x &= -12 &&\justif{les $x$ disparaissent}
\end{calculs}
\explique{$0x=-12$ n'est vrai pour aucune valeur de $x$ : $\res{S=\varnothing}$.}{$0\times x$
vaut toujours $0$, jamais $-12$}
\explique{Pour l'équation d), les $x$ disparaissent et il reste une égalité fausse :
l'équation n'a aucune solution.}{réponse à la question}
\end{exo}

\begin{exo}[Équations avec des parenthèses]
\explique{On développe d'abord chaque membre (un $-$ devant une parenthèse change tous les
signes), on réduit, puis on regroupe les inconnues d'un côté.}{méthode 5 du cours,
exemple b)}
\cas{a)}
\begin{calculs}
3(x-4) &= 2(-x+4) &&\justif{on recopie}\\
3x-12 &= -2x+8 &&\justif{on développe d'abord}\\
3x+2x &= 8+12 &&\\
5x &= 20 &&\\
x &= 4 &&\justif{on divise par $5$}
\end{calculs}
\explique{$\res{S=\ens{4}}$.}{vérification : $3\times 0=0$ et $2\times 0=0$}
\cas{b)}
\begin{calculs}
-2(-a+5) &= 4(3-a)+8 &&\justif{l'inconnue s'appelle $a$}\\
2a-10 &= 12-4a+8 &&\justif{$-2\times(-a)=2a$}\\
2a-10 &= 20-4a &&\justif{on réduit le membre de droite}\\
2a+4a &= 20+10 &&\\
6a &= 30 &&\\
a &= 5 &&
\end{calculs}
\explique{$\res{S=\ens{5}}$.}{vérification : $-2\times 0=0$ et $4\times(-2)+8=0$}
\cas{c)}
\begin{calculs}
m-(7-2m) &= 2m+5 &&\justif{on recopie}\\
m-7+2m &= 2m+5 &&\justif{le $-$ change les signes de $7-2m$}\\
3m-7 &= 2m+5 &&\\
3m-2m &= 5+7 &&\\
m &= 12 &&
\end{calculs}
\explique{$\res{S=\ens{12}}$.}{vérification : $12-(-17)=29$ et $24+5=29$}
\cas{d)}
\begin{calculs}
3(p-2)+4(p-5) &= 0 &&\justif{on recopie}\\
3p-6+4p-20 &= 0 &&\\
7p-26 &= 0 &&\\
7p &= 26 &&\\
p &= \dfrac{26}{7} &&\justif{on garde la fraction : $26\div 7$ ne tombe pas juste}
\end{calculs}
\explique{$\res{S=\ens{\tfrac{26}{7}}}$.}{}
\rem{En c), erreur fréquente : $-(7-2m)=-7-2m$. Le $-$ change \emph{les deux} signes.}
\end{exo}

\begin{exo}[Équations produit]
\explique{Un produit est nul si et seulement si l'un de ses facteurs est nul : on écrit
« facteur $=0$ » pour chaque facteur. Si l'expression n'est pas un produit, on la factorise
d'abord.}{propriété 3 et méthode 5 du cours, exemples c) à e)}
\cas{a)}
\begin{calculs}
(3x-9)(x+5) &= 0 &&\justif{on recopie}\\
3x-9 &= 0 \quad\text{ou}\quad x+5=0 &&\justif{produit nul : un des facteurs est nul}\\
x &= 3 \quad\text{ou}\quad x=-5 &&\justif{$3x=9$ ; $x=-5$}
\end{calculs}
\explique{$\res{S=\ens{-5\,;3}}$.}{solutions rangées dans l'ordre croissant}
\cas{b)}
\begin{calculs}
(2x+10)(3-4x) &= 0 &&\justif{on recopie}\\
2x+10 &= 0 \quad\text{ou}\quad 3-4x=0 &&\justif{produit nul}\\
x &= -5 \quad\text{ou}\quad x=\dfrac{3}{4} &&\justif{$-4x=-3$, donc $x=\frac{-3}{-4}$}
\end{calculs}
\explique{$\res{S=\ens{-5\,;\tfrac{3}{4}}}$.}{}
\cas{c)}
\begin{calculs}
4x(3x-8) &= 0 &&\justif{deux facteurs : $4x$ et $3x-8$}\\
4x &= 0 \quad\text{ou}\quad 3x-8=0 &&\\
x &= 0 \quad\text{ou}\quad x=\dfrac{8}{3} &&
\end{calculs}
\explique{$\res{S=\ens{0\,;\tfrac{8}{3}}}$.}{}
\cas{d)}
\begin{calculs}
(2x+5)^{2} &= 0 &&\justif{on recopie}\\
(2x+5)(2x+5) &= 0 &&\justif{deux facteurs égaux}\\
2x+5 &= 0 &&\justif{une seule équation}\\
x &= -\dfrac{5}{2} &&
\end{calculs}
\explique{$\res{S=\ens{-\tfrac{5}{2}}}$.}{une seule solution}
\cas{e)}
\begin{calculs}
6x^{2}-5x &= 0 &&\justif{ce n'est pas un produit : on factorise}\\
x(6x-5) &= 0 &&\justif{facteur commun $x$}\\
x &= 0 \quad\text{ou}\quad 6x-5=0 &&\justif{produit nul}\\
x &= 0 \quad\text{ou}\quad x=\dfrac{5}{6} &&
\end{calculs}
\explique{$\res{S=\ens{0\,;\tfrac{5}{6}}}$.}{on ne divise jamais par $x$ : on perdrait la
solution $0$}
\cas{f)}
\begin{calculs}
x^{2}-49 &= 0 &&\justif{différence de deux carrés}\\
(x-7)(x+7) &= 0 &&\justif{$a^{2}-b^{2}$, avec $a=x$ et $b=7$}\\
x-7 &= 0 \quad\text{ou}\quad x+7=0 &&\\
x &= 7 \quad\text{ou}\quad x=-7 &&
\end{calculs}
\explique{$\res{S=\ens{-7\,;7}}$.}{autre rédaction : $x^{2}=49$ et $49>0$ (propriété 4)}
\cas{g)}
\begin{calculs}
(x-3)(2x+5)+(x-3)(x-8) &= 0 &&\justif{on recopie}\\
(x-3)\left[(2x+5)+(x-8)\right] &= 0 &&\justif{on factorise par $(x-3)$}\\
(x-3)(3x-3) &= 0 &&\justif{$2x+x=3x$ ; $5-8=-3$}\\
x-3 &= 0 \quad\text{ou}\quad 3x-3=0 &&\justif{produit nul}\\
x &= 3 \quad\text{ou}\quad x=1 &&
\end{calculs}
\explique{$\res{S=\ens{1\,;3}}$.}{}
\cas{h)}
\begin{calculs}
(3x+1)(x+2)-(3x+1)(2x+6) &= 0 &&\justif{on recopie}\\
(3x+1)\left[(x+2)-(2x+6)\right] &= 0 &&\justif{on factorise par $(3x+1)$}\\
(3x+1)(x+2-2x-6) &= 0 &&\justif{le $-$ change les signes de $2x+6$}\\
(3x+1)(-x-4) &= 0 &&\\
3x+1 &= 0 \quad\text{ou}\quad -x-4=0 &&\justif{produit nul}\\
x &= -\dfrac{1}{3} \quad\text{ou}\quad x=-4 &&
\end{calculs}
\explique{$\res{S=\ens{-4\,;-\tfrac{1}{3}}}$.}{$-4<-\frac{1}{3}$ : ordre croissant}
\rem{Erreur fréquente en e) : diviser par $x$ et ne trouver que $\frac{5}{6}$. On factorise
toujours.}
\end{exo}

\begin{exo}[Un premier problème]
\cas{Résolution}
\explique{Soit $c$ le nombre de condensateurs. Il y a deux fois plus de résistances : $2c$.
Il y a $8$ condensateurs de plus que de diodes, donc $8$ diodes de moins que de
condensateurs : $c-8$.}{\textbf{inconnue} : on écrit ce que représente la lettre}
\begin{calculs}
2c+c+(c-8) &= 96 &&\justif{\textbf{équation} : les $96$ composants}\\
4c-8 &= 96 &&\justif{on réduit}\\
4c &= 104 &&\\
c &= 26 &&\justif{on divise par $4$}
\end{calculs}
\explique{$2\times 26=52$ résistances et $26-8=18$ diodes ; $52+26+18=96$.}{\textbf{vérification}}
\explique{Le magasin a reçu \res{$52$ résistances, $26$ condensateurs et $18$ diodes}.}{%
\textbf{conclusion} : une phrase qui répond à la question}
\rem{« $8$ condensateurs de plus que de diodes » : ce sont les diodes qui sont $8$ de moins,
$c-8$ et non $c+8$.}
\end{exo}

\partie{Mise en équation}

\begin{exo}[Deux commandes de câbles]
\cas{a)}
\explique{Le premier club achète $x$ câbles à $4$~€ : il dépense $\res{4x}$ euros.}{prix d'un
câble $\times$ nombre de câbles}
\explique{Le second club achète $15$ câbles de plus : $\res{x+15}$ câbles.}{}
\explique{Il les paie $1{,}50$~€ l'un : il dépense $\res{1{,}5(x+15)}$ euros.}{parenthèses :
chaque câble coûte $1{,}50$~€}
\cas{b)}
\explique{Soit $x$ le nombre de câbles achetés par le premier club. Les deux dépenses sont
égales.}{\textbf{inconnue}, puis \textbf{équation}}
\begin{calculs}
4x &= 1{,}5(x+15) &&\justif{on traduit « la même somme »}\\
4x &= 1{,}5x+22{,}5 &&\justif{$1{,}5\times 15=22{,}5$}\\
4x-1{,}5x &= 22{,}5 &&\\
2{,}5x &= 22{,}5 &&\\
x &= 9 &&\justif{$22{,}5\div 2{,}5=9$}
\end{calculs}
\explique{Premier club : $4\times 9=36$~€ ; second club : $9+15=24$ câbles et
$1{,}5\times 24=36$~€.}{\textbf{vérification}}
\explique{Le premier club a acheté \res{$9$ câbles}.}{\textbf{conclusion}}
\end{exo}

\begin{exo}[Une question d'âges]
\cas{Résolution}
\explique{Soit $x$ le nombre d'années qui passent. La technicienne aura $58+x$ ans, sa
fille $33+x$ ans et son petit-fils $7+x$ ans.}{\textbf{inconnue} : chacun vieillit de $x$ ans}
\begin{calculs}
(33+x)+(7+x) &= 58+x &&\justif{\textbf{équation} : la somme des deux âges, égale au sien}\\
40+2x &= 58+x &&\justif{on réduit}\\
2x-x &= 58-40 &&\\
x &= 18 &&
\end{calculs}
\explique{Dans $18$ ans : la fille aura $51$ ans, le petit-fils $25$ ans ; $51+25=76$, l'âge de
la technicienne à ce moment.}{\textbf{vérification}}
\explique{L'égalité n'arrive que lorsque la technicienne a $76$ ans, après ses $70$ ans :
\res{elle a tort}.}{\textbf{conclusion} : on compare à la condition de l'énoncé}
\rem{Trouver $x$ ne suffit pas : on revient à la question (avant $70$ ans ?). Erreur
fréquente : ne faire vieillir que la technicienne, $33+7=58+x$.}
\end{exo}

\begin{exo}[Trottinettes en libre-service]
\cas{a)}
\begin{calculs}
\text{prix} &= 1+0{,}25\times 12 &&\justif{déverrouillage, puis $12$ minutes}\\
\text{prix} &= 1+3 &&\\
\text{prix} &= 4 &&
\end{calculs}
\explique{Un trajet de $12$ minutes coûte \res{$4$~€}.}{}
\cas{b)}
\explique{Soit $x$ la durée du trajet, en minutes.}{\textbf{inconnue}}
\begin{calculs}
1+0{,}25x &= 6{,}5 &&\justif{\textbf{équation}}\\
0{,}25x &= 5{,}5 &&\justif{on retire $1$ aux deux membres}\\
x &= 22 &&\justif{$5{,}5\div 0{,}25=22$ (diviser par $0{,}25$, c'est multiplier par $4$)}
\end{calculs}
\explique{$1+0{,}25\times 22=1+5{,}5$, soit $6{,}50$~€.}{\textbf{vérification}}
\explique{Le trajet a duré \res{$22$ minutes}.}{\textbf{conclusion}}
\cas{c)}
\explique{Soit $x$ la durée du trajet, en minutes. Second opérateur : $0{,}35x$ euros.}{}
\begin{calculs}
1+0{,}25x &= 0{,}35x &&\justif{même prix chez les deux opérateurs}\\
1 &= 0{,}35x-0{,}25x &&\justif{les $x$ du côté où il y en a le plus}\\
1 &= 0{,}1x &&\\
x &= 10 &&\justif{$1\div 0{,}1=10$}
\end{calculs}
\explique{$1+0{,}25\times 10=3{,}5$ et $0{,}35\times 10=3{,}5$.}{\textbf{vérification}}
\explique{Les deux opérateurs coûtent le même prix pour un trajet de \res{$10$ minutes}
($3{,}50$~€).}{\textbf{conclusion}}
\end{exo}

\partie{Résoudre des inéquations}

\begin{exo}[Inéquations immédiates]
\explique{On isole $x$ comme dans une équation ; quand on divise par un nombre
\textbf{négatif}, le sens de l'inégalité change.}{propriété 5 et méthode 7 du cours}
\cas{a)}
\begin{calculs}
x-5 &\le 0 &&\justif{on recopie}\\
x &\le 5 &&\justif{on ajoute $5$}
\end{calculs}
\explique{$\res{S=\left]-\infty\,;5\right]}$.}{$5$ est compris : inégalité large}
\cas{b)}
\begin{calculs}
x+3 &> 0 &&\\
x &> -3 &&\justif{on retire $3$}
\end{calculs}
\explique{$\res{S=\left]-3\,;+\infty\right[}$.}{$-3$ est exclu : inégalité stricte}
\cas{c)}
\begin{calculs}
4x &\le 20 &&\\
x &\le 5 &&\justif{on divise par $4>0$ : le sens ne change pas}
\end{calculs}
\explique{$\res{S=\left]-\infty\,;5\right]}$.}{}
\cas{d)}
\begin{calculs}
2x+9 &> 0 &&\\
2x &> -9 &&\\
x &> -\dfrac{9}{2} &&\justif{on divise par $2>0$}
\end{calculs}
\explique{$\res{S=\left]-\tfrac{9}{2}\,;+\infty\right[}$.}{}
\cas{e)}
\begin{calculs}
-3x &> 12 &&\\
x &< -4 &&\justif{on divise par $-3<0$ : \textbf{le sens change}}
\end{calculs}
\explique{$\res{S=\left]-\infty\,;-4\right[}$.}{contrôle : $x=-5$ donne $15>12$}
\cas{f)}
\begin{calculs}
-x+7 &\ge 0 &&\\
-x &\ge -7 &&\justif{on retire $7$}\\
x &\le 7 &&\justif{on divise par $-1<0$ : le sens change}
\end{calculs}
\explique{$\res{S=\left]-\infty\,;7\right]}$.}{}
\cas{g)}
\begin{calculs}
5x-2 &< 0 &&\\
5x &< 2 &&\\
x &< \dfrac{2}{5} &&\justif{on divise par $5>0$}
\end{calculs}
\explique{$\res{S=\left]-\infty\,;\tfrac{2}{5}\right[}$.}{}
\cas{h)}
\begin{calculs}
\dfrac{3-4x}{5} &\ge 0 &&\justif{on recopie}\\
3-4x &\ge 0 &&\justif{on multiplie par $5>0$ : le sens ne change pas}\\
-4x &\ge -3 &&\\
x &\le \dfrac{3}{4} &&\justif{on divise par $-4<0$ : \textbf{le sens change}}
\end{calculs}
\explique{$\res{S=\left]-\infty\,;\tfrac{3}{4}\right]}$.}{}
\rem{Crochet fermé pour $\le$ ou $\ge$, ouvert pour $<$ ou $>$ ; toujours ouvert du côté de
l'infini.}
\end{exo}

\begin{exo}[L'inconnue des deux côtés]
\explique{Comme pour une équation : on développe, on regroupe les $x$ d'un côté et les
nombres de l'autre, puis on divise, en regardant le signe du diviseur.}{méthode 7 du cours}
\cas{a)}
\begin{calculs}
4x-5 &< 3-2x &&\justif{on recopie}\\
4x+2x &< 3+5 &&\\
6x &< 8 &&\\
x &< \dfrac{8}{6} &&\justif{on divise par $6>0$}\\
x &< \dfrac{\barre{2}\times 4}{\barre{2}\times 3} &&\justif{on simplifie par $2$}\\
x &< \dfrac{4}{3} &&
\end{calculs}
\explique{$\res{S=\left]-\infty\,;\tfrac{4}{3}\right[}$.}{}
\cas{b)}
\begin{calculs}
3x-2 &\ge 6x+7 &&\justif{on recopie}\\
3x-6x &\ge 7+2 &&\\
-3x &\ge 9 &&\\
x &\le -3 &&\justif{on divise par $-3<0$ : \textbf{le sens change}}
\end{calculs}
\explique{$\res{S=\left]-\infty\,;-3\right]}$.}{contrôle : $x=-4$ donne $-14\ge -17$}
\cas{c)}
\begin{calculs}
3(x-2) &\ge 9-2x &&\justif{on recopie}\\
3x-6 &\ge 9-2x &&\justif{on développe}\\
3x+2x &\ge 9+6 &&\\
5x &\ge 15 &&\\
x &\ge 3 &&\justif{on divise par $5>0$}
\end{calculs}
\explique{$\res{S=\left[3\,;+\infty\right[}$.}{}
\cas{d)}
\begin{calculs}
-2(x-3) &< 4-x &&\justif{on recopie}\\
-2x+6 &< 4-x &&\justif{$-2\times(-3)=+6$}\\
-2x+x &< 4-6 &&\\
-x &< -2 &&\\
x &> 2 &&\justif{on divise par $-1<0$ : le sens change}
\end{calculs}
\explique{$\res{S=\left]2\,;+\infty\right[}$.}{}
\cas{e)}
\begin{calculs}
3(x+1) &< 5+3x &&\justif{on recopie}\\
3x+3 &< 5+3x &&\\
3x-3x &< 5-3 &&\\
0x &< 2 &&\justif{les $x$ disparaissent}
\end{calculs}
\explique{$0<2$ est vrai pour toutes les valeurs de $x$ : $\res{S=\R}$.}{$0\times x$ vaut
toujours $0$}
\cas{f)}
\begin{calculs}
2(x-4) &\ge 2x+1 &&\justif{on recopie}\\
2x-8 &\ge 2x+1 &&\\
2x-2x &\ge 1+8 &&\\
0x &\ge 9 &&\justif{les $x$ disparaissent}
\end{calculs}
\explique{$0\ge 9$ est faux pour toutes les valeurs de $x$ : $\res{S=\varnothing}$.}{}
\explique{Pour e) et f), les $x$ disparaissent : il reste une inégalité toujours vraie
(tous les réels sont solutions) ou toujours fausse (aucune solution).}{réponse à la
question}
\end{exo}

\begin{exo}[Inéquations avec des fractions]
\explique{On multiplie les deux membres par un nombre \textbf{positif} qui supprime les
dénominateurs : le sens ne change pas. Puis on résout comme d'habitude.}{méthode 7 du cours}
\cas{a)}
\begin{calculs}
\dfrac{x}{2}+3 &\le 5 &&\justif{on recopie}\\
\dfrac{x}{2} &\le 2 &&\justif{on retire $3$}\\
x &\le 4 &&\justif{on multiplie par $2>0$}
\end{calculs}
\explique{$\res{S=\left]-\infty\,;4\right]}$.}{}
\cas{b)}
\begin{calculs}
\dfrac{3x-2}{4} &< 4 &&\justif{on recopie}\\
3x-2 &< 16 &&\justif{on multiplie par $4>0$}\\
3x &< 18 &&\\
x &< 6 &&\justif{on divise par $3>0$}
\end{calculs}
\explique{$\res{S=\left]-\infty\,;6\right[}$.}{}
\cas{c)}
\begin{calculs}
\dfrac{x-1}{2}-\dfrac{2-x}{3} &\ge 0 &&\justif{on recopie}\\
3(x-1)-2(2-x) &\ge 0 &&\justif{on multiplie par $6>0$ : $6\times\frac{x-1}{2}=3(x-1)$}\\
3x-3-4+2x &\ge 0 &&\justif{le $-2$ multiplie chaque terme de $(2-x)$}\\
5x-7 &\ge 0 &&\\
5x &\ge 7 &&\\
x &\ge \dfrac{7}{5} &&\justif{on divise par $5>0$}
\end{calculs}
\explique{$\res{S=\left[\tfrac{7}{5}\,;+\infty\right[}$.}{}
\cas{d)}
\begin{calculs}
\dfrac{x}{4}-\dfrac{6-x}{2} &> 3 &&\justif{on recopie}\\
x-2(6-x) &> 12 &&\justif{on multiplie par $4>0$ ; sans oublier le $3$}\\
x-12+2x &> 12 &&\\
3x &> 24 &&\\
x &> 8 &&
\end{calculs}
\explique{$\res{S=\left]8\,;+\infty\right[}$.}{contrôle : $x=10$ donne $2{,}5+2=4{,}5>3$}
\rem{Un numérateur précédé d'un $-$ se met entre parenthèses : $-\frac{2-x}{3}$ devient
$-2(2-x)$, pas $-4-2x$. Et on multiplie \emph{tous} les termes, y compris le $3$ de d).}
\end{exo}

\partie{Tableaux de signes}

\begin{exo}[Signe d'une expression affine]
\explique{On cherche où $ax+b$ s'annule ; après cette valeur, l'expression a le signe de
$a$, avant, le signe contraire.}{propriété 6 du cours}
\cas{1)}
\begin{calculs}
3x-12 &= 0 &&\justif{on cherche où l'expression s'annule}\\
3x &= 12 &&\\
x &= 4 &&
\end{calculs}
\explique{$a=3>0$ : $-$ puis $+$.}{après $4$, le signe de $a$}
\begin{tabsignes}[2.4]{$x$/1,$3x-12$/1}{$-\infty$,$4$,$+\infty$}
  \tkzTabLine{,-,z,+,}
\end{tabsignes}
\cas{2)}
\begin{calculs}
-2x+10 &= 0 &&\\
-2x &= -10 &&\\
x &= 5 &&
\end{calculs}
\explique{$a=-2<0$ : $+$ puis $-$.}{après $5$, le signe de $a$, négatif}
\begin{tabsignes}[2.4]{$x$/1,$-2x+10$/1}{$-\infty$,$5$,$+\infty$}
  \tkzTabLine{,+,z,-,}
\end{tabsignes}
\cas{3)}
\explique{$x+5=0$ pour $x=-5$ ; $a=1>0$ : $-$ puis $+$.}{$x+5$ s'écrit $1x+5$}
\begin{tabsignes}[2.4]{$x$/1,$x+5$/1}{$-\infty$,$-5$,$+\infty$}
  \tkzTabLine{,-,z,+,}
\end{tabsignes}
\cas{4)}
\explique{$7-x=0$ pour $x=7$ ; $a=-1<0$ : $+$ puis $-$.}{$7-x$ s'écrit $-1x+7$}
\begin{tabsignes}[2.4]{$x$/1,$7-x$/1}{$-\infty$,$7$,$+\infty$}
  \tkzTabLine{,+,z,-,}
\end{tabsignes}
\cas{5)}
\begin{calculs}
5x+2 &= 0 &&\\
5x &= -2 &&\\
x &= -\dfrac{2}{5} &&
\end{calculs}
\explique{$a=5>0$ : $-$ puis $+$.}{}
\begin{tabsignes}[2.4]{$x$/1,$5x+2$/1}{$-\infty$,$-\frac{2}{5}$,$+\infty$}
  \tkzTabLine{,-,z,+,}
\end{tabsignes}
\cas{6)}
\begin{calculs}
-4x-6 &= 0 &&\\
-4x &= 6 &&\\
x &= -\dfrac{6}{4} &&\justif{$\frac{6}{-4}$}\\
x &= -\dfrac{3}{2} &&\justif{on simplifie par $2$}
\end{calculs}
\explique{$a=-4<0$ : $+$ puis $-$.}{}
\begin{tabsignes}[2.4]{$x$/1,$-4x-6$/1}{$-\infty$,$-\frac{3}{2}$,$+\infty$}
  \tkzTabLine{,+,z,-,}
\end{tabsignes}
\rem{Erreur fréquente : prendre le signe du nombre $b$. C'est le coefficient $a$ de $x$ qui
compte : $-4x-6$ commence par $+$.}
\end{exo}

\begin{exo}[Signe d'un produit de deux facteurs]
\explique{On cherche où chaque facteur s'annule (valeurs rangées dans l'ordre croissant en
tête du tableau), une ligne par facteur, puis la ligne du produit par la règle des signes,
colonne par colonne.}{méthode 8 du cours}
\cas{1)}
\explique{$x-5=0$ pour $x=5$ ; \quad $x-2=0$ pour $x=2$.}{valeurs qui annulent chaque facteur}
\explique{$x-2$ : $a=1>0$, donc $-$ puis $+$. \\ $x-5$ : $a=1>0$, donc $-$ puis $+$.}{une
ligne par facteur}
\begin{tabsignes}{$x$/1,$x-2$/1,$x-5$/1,$(x-5)(x-2)$/1.2}{$-\infty$,$2$,$5$,$+\infty$}
  \tkzTabLine{,-,z,+,t,+,}
  \tkzTabLine{,-,t,-,z,+,}
  \tkzTabLine{,+,z,-,z,+,}
\end{tabsignes}
\explique{L'expression est positive sur $\left]-\infty\,;2\right[$ et sur
$\left]5\,;+\infty\right[$, négative sur $\left]2\,;5\right[$.}{dernière ligne : règle des
signes}
\cas{2)}
\begin{calculs}
1-3x &= 0 &&\\
-3x &= -1 &&\\
x &= \dfrac{1}{3} &&
\end{calculs}
\explique{$x+4=0$ pour $x=-4$.}{$-4<\frac{1}{3}$ : $-4$ d'abord}
\explique{$x+4$ : $a=1>0$, donc $-$ puis $+$. \\ $1-3x$ : $a=-3<0$, donc $+$ puis $-$.}{}
\begin{tabsignes}{$x$/1,$x+4$/1,$1-3x$/1,$(1-3x)(x+4)$/1.2}{$-\infty$,$-4$,$\frac{1}{3}$,$+\infty$}
  \tkzTabLine{,-,z,+,t,+,}
  \tkzTabLine{,+,t,+,z,-,}
  \tkzTabLine{,-,z,+,z,-,}
\end{tabsignes}
\explique{L'expression est négative sur $\left]-\infty\,;-4\right[$ et sur
$\left]\tfrac{1}{3}\,;+\infty\right[$, positive sur $\left]-4\,;\tfrac{1}{3}\right[$.}{}
\cas{3)}
\explique{$2x+6=0$ pour $x=-3$ ; \quad $5-x=0$ pour $x=5$.}{}
\explique{$2x+6$ : $a=2>0$, donc $-$ puis $+$. \\ $5-x$ : $a=-1<0$, donc $+$ puis $-$.}{}
\begin{tabsignes}{$x$/1,$2x+6$/1,$5-x$/1,$(2x+6)(5-x)$/1.2}{$-\infty$,$-3$,$5$,$+\infty$}
  \tkzTabLine{,-,z,+,t,+,}
  \tkzTabLine{,+,t,+,z,-,}
  \tkzTabLine{,-,z,+,z,-,}
\end{tabsignes}
\explique{L'expression est positive sur $\left]-3\,;5\right[$, négative sur
$\left]-\infty\,;-3\right[$ et sur $\left]5\,;+\infty\right[$.}{}
\cas{4)}
\begin{calculs}
x^{2}-16 &= x^{2}-4^{2} &&\justif{différence de deux carrés}\\
x^{2}-16 &= (x-4)(x+4) &&\justif{on factorise d'abord}
\end{calculs}
\explique{$x+4$ : $-$ puis $+$, nul en $-4$. \\ $x-4$ : $-$ puis $+$, nul en $4$.}{}
\begin{tabsignes}{$x$/1,$x+4$/1,$x-4$/1,$x^{2}-16$/1.2}{$-\infty$,$-4$,$4$,$+\infty$}
  \tkzTabLine{,-,z,+,t,+,}
  \tkzTabLine{,-,t,-,z,+,}
  \tkzTabLine{,+,z,-,z,+,}
\end{tabsignes}
\explique{$x^{2}-16$ est négatif sur $\left]-4\,;4\right[$, positif sur
$\left]-\infty\,;-4\right[$ et sur $\left]4\,;+\infty\right[$.}{}
\rem{On ne dresse pas le tableau de signes d'une somme : on factorise d'abord, pour avoir
des facteurs du premier degré.}
\end{exo}

\begin{exo}[Signe d'un produit de trois facteurs]
\explique{Même méthode avec trois facteurs : trois valeurs en tête du tableau, trois lignes,
puis celle du produit.}{méthode 8 du cours}
\cas{1)}
\explique{Le facteur $x$ est nul pour $x=0$ ; \quad $x-2=0$ pour $x=2$ ; \quad $x+3=0$ pour $x=-3$.}{valeurs
rangées : $-3$ ; $0$ ; $2$}
\explique{Les trois facteurs ont un coefficient $a=1>0$ : chacun est $-$ puis $+$.}{}
\begin{tabsignes}{$x$/1,$x+3$/1,$x$/1,$x-2$/1,$x(x-2)(x+3)$/1.2}{$-\infty$,$-3$,$0$,$2$,$+\infty$}
  \tkzTabLine{,-,z,+,t,+,t,+,}
  \tkzTabLine{,-,t,-,z,+,t,+,}
  \tkzTabLine{,-,t,-,t,-,z,+,}
  \tkzTabLine{,-,z,+,z,-,z,+,}
\end{tabsignes}
\explique{Le produit est négatif sur $\left]-\infty\,;-3\right[$ et sur $\left]0\,;2\right[$,
positif sur $\left]-3\,;0\right[$ et sur $\left]2\,;+\infty\right[$.}{trois $-$ : négatif ;
deux $-$ : positif}
\cas{2)}
\explique{$3x=0$ pour $x=0$ ; \quad $2x-5=0$ pour $x=\dfrac{5}{2}$ ; \quad $x+4=0$ pour
$x=-4$.}{valeurs rangées : $-4$ ; $0$ ; $\frac{5}{2}$}
\explique{$3x$, $2x-5$ et $x+4$ ont des coefficients positifs : chacun est $-$ puis $+$.}{}
\begin{tabsignes}{$x$/1,$x+4$/1,$3x$/1,$2x-5$/1,$3x(2x-5)(x+4)$/1.2}{$-\infty$,$-4$,$0$,$\frac{5}{2}$,$+\infty$}
  \tkzTabLine{,-,z,+,t,+,t,+,}
  \tkzTabLine{,-,t,-,z,+,t,+,}
  \tkzTabLine{,-,t,-,t,-,z,+,}
  \tkzTabLine{,-,z,+,z,-,z,+,}
\end{tabsignes}
\explique{Le produit est négatif sur $\left]-\infty\,;-4\right[$ et sur
$\left]0\,;\tfrac{5}{2}\right[$, positif sur $\left]-4\,;0\right[$ et sur
$\left]\tfrac{5}{2}\,;+\infty\right[$.}{}
\cas{3)}
\begin{calculs}
4-x^{2} &= 2^{2}-x^{2} &&\justif{différence de deux carrés}\\
4-x^{2} &= (2-x)(2+x) &&\justif{$a=2$ et $b=x$}
\end{calculs}
\explique{Le produit s'écrit $(2-x)(2+x)(x-3)$ ; les facteurs s'annulent en $2$, en $-2$ et
en $3$.}{valeurs rangées : $-2$ ; $2$ ; $3$}
\explique{$2+x$ : $-$ puis $+$. \\ $2-x$ : $a=-1<0$, donc $+$ puis $-$. \\ $x-3$ : $-$ puis
$+$.}{}
\begin{tabsignes}{$x$/1,$2+x$/1,$2-x$/1,$x-3$/1,$\left(4-x^{2}\right)(x-3)$/1.2}{$-\infty$,$-2$,$2$,$3$,$+\infty$}
  \tkzTabLine{,-,z,+,t,+,t,+,}
  \tkzTabLine{,+,t,+,z,-,t,-,}
  \tkzTabLine{,-,t,-,t,-,z,+,}
  \tkzTabLine{,+,z,-,z,+,z,-,}
\end{tabsignes}
\explique{Le produit est positif sur $\left]-\infty\,;-2\right[$ et sur $\left]2\,;3\right[$,
négatif sur $\left]-2\,;2\right[$ et sur $\left]3\,;+\infty\right[$.}{}
\rem{Le facteur $2-x$ est positif \emph{avant} $2$ : ici, le produit commence par $+$. On
ne devine pas l'alternance, on fait la règle des signes colonne par colonne.}
\end{exo}

\partie{Inéquations produit}

\begin{exo}[Deux facteurs]
\explique{On dresse le tableau de signes du produit, puis on lit les solutions sur la
dernière ligne : les $+$ pour $>0$, les $-$ pour $<0$, avec les $0$ si l'inégalité est
large.}{méthode 9 du cours}
\cas{a)}
\explique{Le facteur $x$ est nul pour $x=0$ ; \quad $x-3=0$ pour $x=3$.}{valeurs qui annulent les facteurs}
\begin{tabsignes}{$x$/1,$x$/1,$x-3$/1,$x(x-3)$/1.2}{$-\infty$,$0$,$3$,$+\infty$}
  \tkzTabLine{,-,z,+,t,+,}
  \tkzTabLine{,-,t,-,z,+,}
  \tkzTabLine{,+,z,-,z,+,}
\end{tabsignes}
\explique{$\res{S=\left]-\infty\,;0\right]\cup\left[3\,;+\infty\right[}$.}{on lit les $+$ et
les $0$ : bornes comprises}
\cas{b)}
\explique{$x+2=0$ pour $x=-2$ ; \quad $x-6=0$ pour $x=6$.}{}
\begin{tabsignes}{$x$/1,$x+2$/1,$x-6$/1,$(x+2)(x-6)$/1.2}{$-\infty$,$-2$,$6$,$+\infty$}
  \tkzTabLine{,-,z,+,t,+,}
  \tkzTabLine{,-,t,-,z,+,}
  \tkzTabLine{,+,z,-,z,+,}
\end{tabsignes}
\explique{$\res{S=\left]-2\,;6\right[}$.}{on lit les $-$ ; inégalité stricte : bornes exclues}
\cas{c)}
\explique{$5-x=0$ pour $x=5$ ; \quad $3x+9=0$ pour $x=-3$.}{}
\explique{$3x+9$ : $-$ puis $+$. \\ $5-x$ : $a=-1<0$, donc $+$ puis $-$.}{}
\begin{tabsignes}{$x$/1,$3x+9$/1,$5-x$/1,$(5-x)(3x+9)$/1.2}{$-\infty$,$-3$,$5$,$+\infty$}
  \tkzTabLine{,-,z,+,t,+,}
  \tkzTabLine{,+,t,+,z,-,}
  \tkzTabLine{,-,z,+,z,-,}
\end{tabsignes}
\explique{$\res{S=\left[-3\,;5\right]}$.}{on lit les $+$ et les $0$}
\cas{d)}
\explique{$3x-4=0$ pour $x=\dfrac{4}{3}$ ; \quad $1-5x=0$ pour $x=\dfrac{1}{5}$.}{$\frac{1}{5}<\frac{4}{3}$
: $\frac{1}{5}$ d'abord}
\explique{$1-5x$ : $a=-5<0$, donc $+$ puis $-$. \\ $3x-4$ : $-$ puis $+$.}{}
\begin{tabsignes}{$x$/1,$1-5x$/1,$3x-4$/1,$(3x-4)(1-5x)$/1.2}{$-\infty$,$\frac{1}{5}$,$\frac{4}{3}$,$+\infty$}
  \tkzTabLine{,+,z,-,t,-,}
  \tkzTabLine{,-,t,-,z,+,}
  \tkzTabLine{,-,z,+,z,-,}
\end{tabsignes}
\explique{$\res{S=\left]-\infty\,;\tfrac{1}{5}\right[\cup\left]\tfrac{4}{3}\,;+\infty\right[}$.}{on
lit les $-$ ; deux intervalles réunis par $\cup$}
\rem{Inégalité large ($\le$, $\ge$) : les valeurs qui annulent sont solutions, crochets
fermés. Inégalité stricte : crochets ouverts.}
\end{exo}

\begin{exo}[Trois facteurs]
\explique{Trois facteurs : trois valeurs en tête du tableau, trois lignes, puis celle du
produit.}{méthode 9 du cours, exemple c)}
\cas{a)}
\explique{Les facteurs $x$, $x-2$ et $x+4$ s'annulent en $0$, en $2$ et en $-4$ ; tous sont
$-$ puis $+$.}{valeurs rangées : $-4$ ; $0$ ; $2$}
\begin{tabsignes}{$x$/1,$x+4$/1,$x$/1,$x-2$/1,$x(x-2)(x+4)$/1.2}{$-\infty$,$-4$,$0$,$2$,$+\infty$}
  \tkzTabLine{,-,z,+,t,+,t,+,}
  \tkzTabLine{,-,t,-,z,+,t,+,}
  \tkzTabLine{,-,t,-,t,-,z,+,}
  \tkzTabLine{,-,z,+,z,-,z,+,}
\end{tabsignes}
\explique{$\res{S=\left[-4\,;0\right]\cup\left[2\,;+\infty\right[}$.}{on lit les $+$ et les
$0$}
\cas{b)}
\explique{$x+3=0$ pour $x=-3$ ; \quad $2-x=0$ pour $x=2$ ; \quad $2x-1=0$ pour
$x=\dfrac{1}{2}$.}{valeurs rangées : $-3$ ; $\frac{1}{2}$ ; $2$}
\explique{$x+3$ et $2x-1$ : $-$ puis $+$. \\ $2-x$ : $a=-1<0$, donc $+$ puis $-$.}{}
\begin{tabsignes}{$x$/1,$x+3$/1,$2x-1$/1,$2-x$/1,$(x+3)(2-x)(2x-1)$/1.2}{$-\infty$,$-3$,$\frac{1}{2}$,$2$,$+\infty$}
  \tkzTabLine{,-,z,+,t,+,t,+,}
  \tkzTabLine{,-,t,-,z,+,t,+,}
  \tkzTabLine{,+,t,+,t,+,z,-,}
  \tkzTabLine{,+,z,-,z,+,z,-,}
\end{tabsignes}
\explique{$\res{S=\left]-3\,;\tfrac{1}{2}\right[\cup\left]2\,;+\infty\right[}$.}{on lit les
$-$ ; bornes exclues}
\cas{c)}
\explique{$2x=0$ pour $x=0$ ; \quad $3x-4=0$ pour $x=\dfrac{4}{3}$ ; \quad $x+6=0$ pour
$x=-6$.}{valeurs rangées : $-6$ ; $0$ ; $\frac{4}{3}$}
\explique{Les trois facteurs ont un coefficient positif : $-$ puis $+$.}{}
\begin{tabsignes}{$x$/1,$x+6$/1,$2x$/1,$3x-4$/1,$2x(3x-4)(x+6)$/1.2}{$-\infty$,$-6$,$0$,$\frac{4}{3}$,$+\infty$}
  \tkzTabLine{,-,z,+,t,+,t,+,}
  \tkzTabLine{,-,t,-,z,+,t,+,}
  \tkzTabLine{,-,t,-,t,-,z,+,}
  \tkzTabLine{,-,z,+,z,-,z,+,}
\end{tabsignes}
\explique{$\res{S=\left]-\infty\,;-6\right]\cup\left[0\,;\tfrac{4}{3}\right]}$.}{on lit les
$-$ et les $0$}
\cas{d)}
\begin{calculs}
x\left(x^{2}-16\right) &> 0 &&\justif{on recopie}\\
x(x-4)(x+4) &> 0 &&\justif{$x^{2}-16=(x-4)(x+4)$}
\end{calculs}
\explique{Les facteurs s'annulent en $-4$, en $0$ et en $4$ ; tous sont $-$ puis $+$.}{}
\begin{tabsignes}{$x$/1,$x+4$/1,$x$/1,$x-4$/1,$x(x-4)(x+4)$/1.2}{$-\infty$,$-4$,$0$,$4$,$+\infty$}
  \tkzTabLine{,-,z,+,t,+,t,+,}
  \tkzTabLine{,-,t,-,z,+,t,+,}
  \tkzTabLine{,-,t,-,t,-,z,+,}
  \tkzTabLine{,-,z,+,z,-,z,+,}
\end{tabsignes}
\explique{$\res{S=\left]-4\,;0\right[\cup\left]4\,;+\infty\right[}$.}{on lit les $+$}
\rem{En b), le facteur $2-x$ est positif \emph{avant} $2$ : le produit commence par $+$.}
\end{exo}

\begin{exo}[Factoriser avant de résoudre]
\explique{L'expression n'est pas un produit : on la factorise d'abord (facteur commun ou
$a^{2}-b^{2}$), puis on fait le tableau de signes.}{méthode 9 du cours, exemple b)}
\cas{a)}
\begin{calculs}
x^{2}-36 &< 0 &&\justif{on recopie}\\
(x-6)(x+6) &< 0 &&\justif{on factorise : $a^{2}-b^{2}$, avec $a=x$ et $b=6$}
\end{calculs}
\begin{tabsignes}{$x$/1,$x+6$/1,$x-6$/1,$(x-6)(x+6)$/1.2}{$-\infty$,$-6$,$6$,$+\infty$}
  \tkzTabLine{,-,z,+,t,+,}
  \tkzTabLine{,-,t,-,z,+,}
  \tkzTabLine{,+,z,-,z,+,}
\end{tabsignes}
\explique{$\res{S=\left]-6\,;6\right[}$.}{on lit les $-$ ; bornes exclues}
\cas{b)}
\begin{calculs}
x^{2}-8x &\le 0 &&\justif{on recopie}\\
x(x-8) &\le 0 &&\justif{on factorise par $x$}
\end{calculs}
\begin{tabsignes}{$x$/1,$x$/1,$x-8$/1,$x(x-8)$/1.2}{$-\infty$,$0$,$8$,$+\infty$}
  \tkzTabLine{,-,z,+,t,+,}
  \tkzTabLine{,-,t,-,z,+,}
  \tkzTabLine{,+,z,-,z,+,}
\end{tabsignes}
\explique{$\res{S=\left[0\,;8\right]}$.}{on lit les $-$ et les $0$}
\cas{c)}
\begin{calculs}
\left(9x^{2}-4\right)(x+2) &> 0 &&\justif{on recopie}\\
(3x-2)(3x+2)(x+2) &> 0 &&\justif{$9x^{2}-4=(3x)^{2}-2^{2}$}
\end{calculs}
\explique{$3x-2=0$ pour $x=\dfrac{2}{3}$ ; \quad $3x+2=0$ pour $x=-\dfrac{2}{3}$ ; \quad
$x+2=0$ pour $x=-2$.}{valeurs rangées : $-2$ ; $-\frac{2}{3}$ ; $\frac{2}{3}$}
\begin{tabsignes}{$x$/1,$x+2$/1,$3x+2$/1,$3x-2$/1,$\left(9x^{2}-4\right)(x+2)$/1.2}{$-\infty$,$-2$,$-\frac{2}{3}$,$\frac{2}{3}$,$+\infty$}
  \tkzTabLine{,-,z,+,t,+,t,+,}
  \tkzTabLine{,-,t,-,z,+,t,+,}
  \tkzTabLine{,-,t,-,t,-,z,+,}
  \tkzTabLine{,-,z,+,z,-,z,+,}
\end{tabsignes}
\explique{$\res{S=\left]-2\,;-\tfrac{2}{3}\right[\cup\left]\tfrac{2}{3}\,;+\infty\right[}$.}{on
lit les $+$}
\cas{d)}
\begin{calculs}
(x+2)^{2}-16 &> 0 &&\justif{$a^{2}-b^{2}$, avec $a=x+2$ et $b=4$}\\
\left[(x+2)-4\right]\left[(x+2)+4\right] &> 0 &&\\
(x-2)(x+6) &> 0 &&\justif{on réduit chaque crochet}
\end{calculs}
\begin{tabsignes}{$x$/1,$x+6$/1,$x-2$/1,$(x-2)(x+6)$/1.2}{$-\infty$,$-6$,$2$,$+\infty$}
  \tkzTabLine{,-,z,+,t,+,}
  \tkzTabLine{,-,t,-,z,+,}
  \tkzTabLine{,+,z,-,z,+,}
\end{tabsignes}
\explique{$\res{S=\left]-\infty\,;-6\right[\cup\left]2\,;+\infty\right[}$.}{on lit les $+$}
\rem{Erreur fréquente en a) : $x^{2}<36$, donc $x<6$. On oublie alors que $-7$, par
exemple, n'est pas solution ($49>36$) : le tableau de signes évite cette erreur.}
\end{exo}

\partie{Inéquations quotient}

\begin{exo}[Premières inéquations quotient]
\cas{1)}
\explique{$x+5=0$ pour $x=-5$ : la valeur interdite de $\dfrac{4-x}{x+5}$ est $\res{-5}$.}{on
cherche où le dénominateur s'annule}
\explique{$2-x=0$ pour $x=2$ : la valeur interdite de $\dfrac{7-2x}{2-x}$ est $\res{2}$.}{}
\explique{$3x-9=0$ pour $x=3$ : la valeur interdite de $\dfrac{x+2}{3x-9}$ est $\res{3}$.}{$3x=9$}
\cas{2) a)}
\explique{Valeur interdite : $-5$. \\ Numérateur : $4-x=0$ pour $x=4$.}{le dénominateur
d'abord (définition 5 et méthode 10 du cours)}
\begin{tabsignes}[2.8]{$x$/1,$4-x$/1,$x+5$/1,$\frac{4-x}{x+5}$/1.4}{$-\infty$,$-5$,$4$,$+\infty$}
  \tkzTabLine{,+,t,+,z,-,}
  \tkzTabLine{,-,z,+,t,+,}
  \tkzTabLine{,-,d,+,z,-,}
\end{tabsignes}
\explique{$\res{S=\left]-5\,;4\right[}$.}{on lit les $+$ ; double barre en $-5$}
\cas{b)}
\explique{Valeur interdite : $3$. \\ Numérateur : $x+2=0$ pour $x=-2$.}{}
\begin{tabsignes}[2.8]{$x$/1,$x+2$/1,$3x-9$/1,$\frac{x+2}{3x-9}$/1.4}{$-\infty$,$-2$,$3$,$+\infty$}
  \tkzTabLine{,-,z,+,t,+,}
  \tkzTabLine{,-,t,-,z,+,}
  \tkzTabLine{,+,z,-,d,+,}
\end{tabsignes}
\explique{$\res{S=\left[-2\,;3\right[}$.}{$-2$ compris (le quotient s'y annule) ; $3$ exclu
(interdit)}
\cas{c)}
\explique{Valeur interdite : $2$. \\ Numérateur : $7-2x=0$ pour $x=\dfrac{7}{2}$.}{$-2x=-7$}
\explique{$7-2x$ et $2-x$ ont un coefficient négatif : $+$ puis $-$.}{}
\begin{tabsignes}[2.8]{$x$/1,$2-x$/1,$7-2x$/1,$\frac{7-2x}{2-x}$/1.4}{$-\infty$,$2$,$\frac{7}{2}$,$+\infty$}
  \tkzTabLine{,+,z,-,t,-,}
  \tkzTabLine{,+,t,+,z,-,}
  \tkzTabLine{,+,d,-,z,+,}
\end{tabsignes}
\explique{$\res{S=\left]-\infty\,;2\right[\cup\left[\tfrac{7}{2}\,;+\infty\right[}$.}{$2$
exclu (interdit) ; $\frac{7}{2}$ compris}
\cas{d)}
\explique{Valeur interdite : $4-x=0$ pour $x=4$. \\ Numérateur : $3x+2=0$ pour
$x=-\dfrac{2}{3}$.}{}
\begin{tabsignes}[2.8]{$x$/1,$3x+2$/1,$4-x$/1,$\frac{3x+2}{4-x}$/1.4}{$-\infty$,$-\frac{2}{3}$,$4$,$+\infty$}
  \tkzTabLine{,-,z,+,t,+,}
  \tkzTabLine{,+,t,+,z,-,}
  \tkzTabLine{,-,z,+,d,-,}
\end{tabsignes}
\explique{$\res{S=\left]-\infty\,;-\tfrac{2}{3}\right[\cup\left]4\,;+\infty\right[}$.}{on lit
les $-$ ; bornes exclues}
\rem{Même avec $\le$ ou $\ge$, la valeur interdite a toujours un crochet ouvert.}
\end{exo}

\begin{exo}[Avec une factorisation]
\explique{On factorise le numérateur ou le dénominateur s'il le faut ; les valeurs interdites
annulent le dénominateur ; puis une ligne par facteur.}{méthode 10 du cours}
\cas{a)}
\explique{Valeur interdite : $5-2x=0$ pour $x=\dfrac{5}{2}$. \\ Numérateur : $x(x-2)$ est nul
en $0$ et en $2$.}{valeurs rangées : $0$ ; $2$ ; $\frac{5}{2}$}
\begin{tabsignes}[2.8]{$x$/1,$x$/1,$x-2$/1,$5-2x$/1,$\frac{x(x-2)}{5-2x}$/1.4}{$-\infty$,$0$,$2$,$\frac{5}{2}$,$+\infty$}
  \tkzTabLine{,-,z,+,t,+,t,+,}
  \tkzTabLine{,-,t,-,z,+,t,+,}
  \tkzTabLine{,+,t,+,t,+,z,-,}
  \tkzTabLine{,+,z,-,z,+,d,-,}
\end{tabsignes}
\explique{$\res{S=\left[0\,;2\right]\cup\left]\tfrac{5}{2}\,;+\infty\right[}$.}{$0$ et $2$
compris ; $\frac{5}{2}$ exclu}
\cas{b)}
\begin{calculs}
\dfrac{x^{2}-16}{2-x} &> 0 &&\justif{on recopie}\\
\dfrac{(x-4)(x+4)}{2-x} &> 0 &&\justif{on factorise le numérateur}
\end{calculs}
\explique{Valeur interdite : $2$. \\ Numérateur nul en $-4$ et en $4$.}{valeurs rangées :
$-4$ ; $2$ ; $4$}
\begin{tabsignes}[2.8]{$x$/1,$x+4$/1,$x-4$/1,$2-x$/1,$\frac{x^{2}-16}{2-x}$/1.4}{$-\infty$,$-4$,$2$,$4$,$+\infty$}
  \tkzTabLine{,-,z,+,t,+,t,+,}
  \tkzTabLine{,-,t,-,t,-,z,+,}
  \tkzTabLine{,+,t,+,z,-,t,-,}
  \tkzTabLine{,+,z,-,d,+,z,-,}
\end{tabsignes}
\explique{$\res{S=\left]-\infty\,;-4\right[\cup\left]2\,;4\right[}$.}{on lit les $+$}
\cas{c)}
\begin{calculs}
\dfrac{6-3x}{x^{2}-9} &\le 0 &&\justif{on recopie}\\
\dfrac{6-3x}{(x-3)(x+3)} &\le 0 &&\justif{on factorise le dénominateur}
\end{calculs}
\explique{Valeurs interdites : $-3$ et $3$. \\ Numérateur : $6-3x=0$ pour $x=2$.}{deux
valeurs interdites, deux doubles barres}
\begin{tabsignes}[2.8]{$x$/1,$x+3$/1,$6-3x$/1,$x-3$/1,$\frac{6-3x}{x^{2}-9}$/1.4}{$-\infty$,$-3$,$2$,$3$,$+\infty$}
  \tkzTabLine{,-,z,+,t,+,t,+,}
  \tkzTabLine{,+,t,+,z,-,t,-,}
  \tkzTabLine{,-,t,-,t,-,z,+,}
  \tkzTabLine{,+,d,-,z,+,d,-,}
\end{tabsignes}
\explique{$\res{S=\left]-3\,;2\right]\cup\left]3\,;+\infty\right[}$.}{$-3$ et $3$ exclus ;
$2$ compris}
\end{exo}

\begin{exo}[Tout ramener d'un même côté]
\explique{On passe tout dans le membre de gauche, on réduit au même dénominateur, puis on
dresse le tableau de signes du quotient obtenu.}{méthode 10 du cours, exemple b)}
\cas{a)}
\explique{Valeur interdite : $-3$.}{}
\begin{calculs}
\dfrac{3x+1}{x+3} &\le 2 &&\justif{on recopie}\\
\dfrac{3x+1}{x+3}-2 &\le 0 &&\justif{tout dans le membre de gauche}\\
\dfrac{3x+1-2(x+3)}{x+3} &\le 0 &&\justif{même dénominateur : $2=\frac{2(x+3)}{x+3}$}\\
\dfrac{x-5}{x+3} &\le 0 &&\justif{$3x+1-2x-6=x-5$}
\end{calculs}
\begin{tabsignes}[2.8]{$x$/1,$x+3$/1,$x-5$/1,$\frac{x-5}{x+3}$/1.4}{$-\infty$,$-3$,$5$,$+\infty$}
  \tkzTabLine{,-,z,+,t,+,}
  \tkzTabLine{,-,t,-,z,+,}
  \tkzTabLine{,+,d,-,z,+,}
\end{tabsignes}
\explique{$\res{S=\left]-3\,;5\right]}$.}{$-3$ exclu (interdit) ; $5$ compris}
\cas{b)}
\explique{Valeur interdite : $2$.}{}
\begin{calculs}
\dfrac{2-5x}{2-x} &\ge 3 &&\justif{on recopie}\\
\dfrac{2-5x}{2-x}-3 &\ge 0 &&\justif{tout dans le membre de gauche}\\
\dfrac{2-5x-3(2-x)}{2-x} &\ge 0 &&\justif{même dénominateur : $3=\frac{3(2-x)}{2-x}$}\\
\dfrac{-2x-4}{2-x} &\ge 0 &&\justif{$2-5x-6+3x=-2x-4$}
\end{calculs}
\explique{$-2x-4=0$ pour $x=-2$ ; $-2x-4$ et $2-x$ ont un coefficient négatif : $+$ puis
$-$.}{}
\begin{tabsignes}[2.8]{$x$/1,$-2x-4$/1,$2-x$/1,$\frac{-2x-4}{2-x}$/1.4}{$-\infty$,$-2$,$2$,$+\infty$}
  \tkzTabLine{,+,z,-,t,-,}
  \tkzTabLine{,+,t,+,z,-,}
  \tkzTabLine{,+,z,-,d,+,}
\end{tabsignes}
\explique{$\res{S=\left]-\infty\,;-2\right]\cup\left]2\,;+\infty\right[}$.}{$-2$ compris ;
$2$ exclu}
\cas{c)}
\explique{Valeur interdite : $3-2x=0$ pour $x=\dfrac{3}{2}$.}{}
\begin{calculs}
\dfrac{x+4}{3-2x} &< 1 &&\justif{on recopie}\\
\dfrac{x+4}{3-2x}-1 &< 0 &&\justif{tout dans le membre de gauche}\\
\dfrac{x+4-(3-2x)}{3-2x} &< 0 &&\justif{même dénominateur : $1=\frac{3-2x}{3-2x}$}\\
\dfrac{3x+1}{3-2x} &< 0 &&\justif{le $-$ change les signes de $3-2x$}
\end{calculs}
\explique{$3x+1=0$ pour $x=-\dfrac{1}{3}$.}{$-\frac{1}{3}<\frac{3}{2}$}
\begin{tabsignes}[2.8]{$x$/1,$3x+1$/1,$3-2x$/1,$\frac{3x+1}{3-2x}$/1.4}{$-\infty$,$-\frac{1}{3}$,$\frac{3}{2}$,$+\infty$}
  \tkzTabLine{,-,z,+,t,+,}
  \tkzTabLine{,+,t,+,z,-,}
  \tkzTabLine{,-,z,+,d,-,}
\end{tabsignes}
\explique{$\res{S=\left]-\infty\,;-\tfrac{1}{3}\right[\cup\left]\tfrac{3}{2}\,;+\infty\right[}$.}{on
lit les $-$ ; bornes exclues}
\rem{Multiplier par $x+3$ serait faux : son signe dépend de $x$, on ne saurait pas s'il faut
changer le sens de l'inégalité.}
\end{exo}

\partie{Problèmes avec des inéquations}

\begin{exo}[Seuil de rentabilité]
\cas{1)}
\begin{calculs}
C(25) &= 15\times 25+480 &&\justif{on remplace $x$ par $25$}\\
C(25) &= 855 &&\justif{$375+480$}\\
R(25) &= 31\times 25 &&\\
R(25) &= 775 &&
\end{calculs}
\explique{$R(25)<C(25)$ : la recette ne couvre pas le coût, l'atelier \res{ne fait pas de
bénéfice} (il perd $80$~€).}{bénéfice : recette supérieure au coût}
\cas{2)}
\explique{L'atelier fait un bénéfice lorsque $R(x)>C(x)$.}{on traduit l'énoncé par une
inéquation}
\begin{calculs}
31x &> 15x+480 &&\justif{on recopie}\\
31x-15x &> 480 &&\\
16x &> 480 &&\\
x &> 30 &&\justif{on divise par $16>0$}
\end{calculs}
\explique{$R(31)=961$ et $C(31)=945$ : bénéfice de $16$~€.}{vérification pour le premier
entier qui convient}
\explique{L'atelier fait un bénéfice à partir de \res{$31$ modules par jour}.}{$x$ est un
nombre entier de modules}
\cas{3)}
\begin{calculs}
\dfrac{15x+480}{x} &\le 21 &&\justif{on traduit l'énoncé}\\
\dfrac{15x+480-21x}{x} &\le 0 &&\justif{tout à gauche, même dénominateur}\\
\dfrac{480-6x}{x} &\le 0 &&\\
480-6x &\le 0 &&\justif{$x>0$ : le quotient a le signe du numérateur}\\
-6x &\le -480 &&\\
x &\ge 80 &&\justif{on divise par $-6<0$ : le sens change}
\end{calculs}
\explique{$C(80)=1\,680$ et $1\,680\div 80=21$ : le coût moyen vaut exactement $21$~€.}{vérification}
\explique{Il faut assembler \res{au moins $80$ modules par jour}.}{conclusion dans le contexte}
\rem{Ici, inutile de faire un tableau de signes : l'énoncé donne $x>0$, donc le dénominateur
est positif.}
\end{exo}

\begin{exo}[Louer un utilitaire électrique]
\cas{1)}
\begin{calculs}
\text{prix A} &= 0{,}5\times 60 &&\justif{pour $60$~km : $0{,}50$~€ par kilomètre}\\
\text{prix A} &= 30 &&\\
\text{prix B} &= 36+0{,}2\times 60 &&\justif{forfait, puis les kilomètres}\\
\text{prix B} &= 48 &&
\end{calculs}
\explique{Pour $60$~km, le loueur \res{A est le moins cher} ($30$~€ contre $48$~€).}{}
\begin{calculs}
\text{prix A} &= 0{,}5\times 150 &&\justif{pour $150$~km}\\
\text{prix A} &= 75 &&\\
\text{prix B} &= 36+0{,}2\times 150 &&\\
\text{prix B} &= 66 &&
\end{calculs}
\explique{Pour $150$~km, le loueur \res{B est le moins cher} ($66$~€ contre $75$~€).}{le
forfait de B est rattrapé sur les longues distances}
\cas{2)}
\explique{$\res{A(x)=0{,}5x}$ \quad et \quad $\res{B(x)=36+0{,}2x}$.}{prix par kilomètre
$\times$ distance, plus le forfait}
\cas{3)}
\begin{calculs}
36+0{,}2x &< 0{,}5x &&\justif{$B(x)<A(x)$}\\
36 &< 0{,}5x-0{,}2x &&\justif{les $x$ du côté où il y en a le plus}\\
36 &< 0{,}3x &&\\
120 &< x &&\justif{on divise par $0{,}3>0$ : $36\div 0{,}3=120$}
\end{calculs}
\explique{Avec $0\le x\le 300$ : $\res{S=\left]120\,;300\right]}$.}{}
\explique{Le loueur B est moins cher que le loueur A pour plus de $120$~km ; à $120$~km, les
deux coûtent $60$~€.}{interprétation}
\cas{4)}
\begin{calculs}
36+0{,}2x &\le 80 &&\justif{$B(x)\le 80$}\\
0{,}2x &\le 44 &&\\
x &\le 220 &&\justif{on divise par $0{,}2>0$ : $44\div 0{,}2=220$}
\end{calculs}
\explique{L'entreprise peut parcourir \res{$220$~km au plus}.}{vérification : $36+44=80$}
\cas{5)}
\explique{Chez C, le prix est $C(x)=60+0{,}1x$. Il faut $C(x)<A(x)$ \textbf{et}
$C(x)<B(x)$.}{« moins cher que les deux autres »}
\begin{calculs}
60+0{,}1x &< 0{,}5x &&\justif{$C(x)<A(x)$}\\
60 &< 0{,}4x &&\\
150 &< x &&\justif{$60\div 0{,}4=150$}
\end{calculs}
\begin{calculs}
60+0{,}1x &< 36+0{,}2x &&\justif{$C(x)<B(x)$}\\
60-36 &< 0{,}2x-0{,}1x &&\\
24 &< 0{,}1x &&\\
240 &< x &&\justif{$24\div 0{,}1=240$}
\end{calculs}
\explique{Les deux conditions ensemble : $x>240$, et $x\le 300$. Le loueur C est strictement
le moins cher pour une distance \res{comprise entre $240$~km (exclu) et $300$~km}, soit
$x\in\left]240\,;300\right]$.}{on garde la condition la plus exigeante}
\rem{Contrôle avec $x=250$ : A coûte $125$~€, B $86$~€ et C $85$~€.}
\end{exo}

\begin{exo}[Bénéfice d'un fabricant de trottinettes]
\cas{1)}
\begin{calculs}
B(10) &= -2\times 10^{2}+240\times 10-4\,000 &&\justif{on remplace $x$ par $10$}\\
B(10) &= -200+2\,400-4\,000 &&\justif{le carré d'abord : $10^{2}=100$}\\
B(10) &= -1\,800 &&\\
B(50) &= -2\times 50^{2}+240\times 50-4\,000 &&\\
B(50) &= -5\,000+12\,000-4\,000 &&\justif{$50^{2}=2\,500$}\\
B(50) &= 3\,000 &&
\end{calculs}
\explique{Pour $10$ trottinettes, $B(10)<0$ : \res{perte de $1\,800$~€}. Pour $50$
trottinettes, \res{bénéfice de $3\,000$~€}.}{bénéfice négatif : c'est une perte}
\cas{2)}
\begin{calculs}
(2x-40)(100-x) &= 200x-2x^{2}-4\,000+40x &&\justif{double distributivité ; $-40\times(-x)=40x$}\\
(2x-40)(100-x) &= -2x^{2}+240x-4\,000 &&\justif{on réduit et on ordonne}
\end{calculs}
\explique{On retrouve l'expression de $B(x)$ : $\res{B(x)=(2x-40)(100-x)}$.}{forme
factorisée du bénéfice}
\cas{3)}
\explique{$2x-40=0$ pour $x=20$ ; \quad $100-x=0$ pour $x=100$.}{valeurs entre $0$ et $120$}
\explique{$2x-40$ : $-$ puis $+$. \\ $100-x$ : $a=-1<0$, donc $+$ puis $-$.}{}
\begin{tabsignes}{$x$/1,$2x-40$/1,$100-x$/1,$B(x)$/1}{$0$,$20$,$100$,$120$}
  \tkzTabLine{,-,z,+,t,+,}
  \tkzTabLine{,+,t,+,z,-,}
  \tkzTabLine{,-,z,+,z,-,}
\end{tabsignes}
\explique{$\res{S=\left]20\,;100\right[}$.}{on lit les $+$ ; bornes exclues}
\cas{4)}
\explique{$x$ est un nombre entier de trottinettes : l'atelier fait un bénéfice s'il en
fabrique \res{de $21$ à $99$ par semaine}.}{$20$ et $100$ donnent un bénéfice nul}
\cas{5)}
\begin{calculs}
-2(x-40)(x-80) &= -2\left(x^{2}-80x-40x+3\,200\right) &&\justif{on développe d'abord le produit}\\
-2(x-40)(x-80) &= -2\left(x^{2}-120x+3\,200\right) &&\\
-2(x-40)(x-80) &= -2x^{2}+240x-6\,400 &&\justif{$-2$ multiplie chaque terme}\\
B(x)-2\,400 &= -2x^{2}+240x-4\,000-2\,400 &&\\
B(x)-2\,400 &= -2x^{2}+240x-6\,400 &&\justif{on retrouve la même expression}
\end{calculs}
\explique{Donc $B(x)-2\,400=-2(x-40)(x-80)$.}{égalité vérifiée}
\begin{calculs}
B(x) &\ge 2\,400 &&\justif{bénéfice d'au moins $2\,400$~€}\\
B(x)-2\,400 &\ge 0 &&\justif{tout dans le membre de gauche}\\
-2(x-40)(x-80) &\ge 0 &&\\
(x-40)(x-80) &\le 0 &&\justif{on divise par $-2<0$ : le sens change}
\end{calculs}
\begin{tabsignes}{$x$/1,$x-40$/1,$x-80$/1,$(x-40)(x-80)$/1.2}{$0$,$40$,$80$,$120$}
  \tkzTabLine{,-,z,+,t,+,}
  \tkzTabLine{,-,t,-,z,+,}
  \tkzTabLine{,+,z,-,z,+,}
\end{tabsignes}
\explique{$S=\left[40\,;80\right]$ : l'atelier doit fabriquer \res{de $40$ à $80$
trottinettes par semaine}.}{on lit les $-$ et les $0$ ; vérification : $B(40)=2\,400$}
\rem{Le tableau est fait sur $\intff{0}{120}$ : la première colonne commence à $0$, pas à
$-\infty$.}
\end{exo}

\end{document}
